\documentclass[11pt,a4paper]{article}

\def\nyear {2025}
\def\nterm {Winter}
\def\nlecturer {Emanuel Milman}
\def\ncourse {Real Functions}

\makeatletter

% packages
\usepackage{amssymb,amsfonts,amsmath,calc,tikz,pgfplots,geometry,mathtools}
\usepackage{color}   % May be necessary if you want to color links
\usepackage[svgnames]{xcolor}
\usepackage[hidelinks]{hyperref}
\usepackage{forest}
\usepackage{commath} % ew
\usepackage{esvect} % \vv makes vectors
\usepackage{centernot} % for the comparison
\usepackage{esdiff}
\usepackage{esint}
\usepackage{amsthm}
\usepackage{fancyhdr}
\usepackage{bm}
\usepackage{witharrows}
\usepackage{bookmark}
\usepackage{tikz-cd}
\usepackage{quiver}
\usepackage{bbm}
\usepackage{textcomp}
\usepackage{float}
\usepackage{gensymb}
\usepackage{cleveref}
\usepackage{environ}
\usepackage{comment}
\usepackage{cancel}
\usepackage{xparse}

% Inevitable cs
\usepackage{listings}
\usepackage[]{algorithm2e}

% tikz libraries
\usetikzlibrary{positioning}
\usetikzlibrary{matrix}
\usetikzlibrary{arrows}
\usetikzlibrary{bending}
\usetikzlibrary{arrows.meta}
\usetikzlibrary{decorations.markings}
\usetikzlibrary{decorations.pathreplacing}
\usetikzlibrary{decorations.markings}
\usetikzlibrary{patterns}
\usetikzlibrary{backgrounds}

% Page style setup
\pagestyle{fancy}
\fancyhfoffset{0pt}
\renewcommand{\headwidth}{\textwidth}
\geometry{margin=1in}
\pgfplotsset{compat=1.18}
\setlength{\headheight}{14.6pt}
\addtolength{\topmargin}{-1.6pt}
\hypersetup{
    colorlinks=false,
    linktoc=section,
    linkcolor=black,
}

%% maketitle setup
\ifx \nauthor\undefined
  \def\nauthor{Yeheli Fomberg}
\else
\fi

\ifx \nemail\undefined
  \def\nemail{\href{mailto:yp@campus.technion.ac.il}
  {\texttt{<yp@campus.technion.ac.il>}}}
\else
\fi

\ifx \ncoursehead \undefined
  \def\ncoursehead{\ncourse}
\fi

\lhead{\emph{\nouppercase{\leftmark}}}
\ifx \nextra \undefined
  \rhead{
    \ifnum\thepage=1
    \else
      \ncoursehead
    \fi}
\else
  \rhead{
    \ifnum\thepage=1
    \else
      \ncoursehead \ (\nextra)
    \fi}
\fi

\def\notes{notes}
\def\problems{problems}

\ifx \ntype \undefined
  \def\ntype{notes}
\else
\fi

\ifx \ntype \notes
  \let\@real@maketitle\maketitle
  \renewcommand{\maketitle}{\@real@maketitle\begin{center}
  \begin{minipage}[c]{0.9\textwidth}\centering\footnotesize
  These notes are not endorsed by the lecturers.
  I have revised them outside lectures to incorporate supplementary
  explanations, clarifications, and material for fun.
  While I have strived for accuracy, any errors or misinterpretations 
  are most likely mine.
  \end{minipage}\end{center}}
\else
  \ifx \ntype \problems
    \let\@real@maketitle\maketitle
    \renewcommand{\maketitle}{\@real@maketitle\begin{center}
    \begin{minipage}[c]{0.9\textwidth}\centering\footnotesize
    These problems are mostly based on the problem sets I was given
    when taking the course.
    
    While I have strived for accuracy and completeness,
    some of the problems would not have been here without the help of
    friends.
    \end{minipage}\end{center}}
  \else
  \fi
\fi


% theorem environments
\theoremstyle{definition}
\newtheorem{definition}{Definition}[section]
\newtheorem{remark}{Remark}[section]
\newtheorem{example}{Example}[section]
\newtheorem{exercise}{Exercise}[section]
\newtheorem{paradox}{Paradox}[section]
\newtheorem{entry}{Entry}[section]
\newtheorem*{solution}{Solution}
\newtheorem*{question}{Question}
\newtheorem*{answer}{Answer}
\newtheorem*{problem}{Problem}
\theoremstyle{plain}
\newtheorem{theorem}{Theorem}[section]
\newtheorem{proposition}[theorem]{Proposition}
\newtheorem{lemma}[theorem]{Lemma}
\newtheorem{corollary}[theorem]{Corollary}
\newenvironment{proofof}[1]{%
    \renewcommand{\proofname}{Proof of #1}%
    \begin{proof}
}{%
    \end{proof}
}
% tikz customization
\pgfarrowsdeclarecombine{twolatex'}{twolatex'}{latex'}{latex'}{latex'}{latex'}
\tikzset{->/.style = {decoration={markings,
                                  mark=at position 0.999
                                  with {\arrow[scale=2]{latex'}}},
                      postaction={decorate}}}
\tikzset{<-/.style = {decoration={markings,
                                  mark=at position 0.001 with {\arrowreversed[scale=2]{latex'}}},
                      postaction={decorate}}}
\tikzset{-->/.style = {decoration={markings,
                                  mark=at position 0.999
                                  with {\arrow[scale=2]{latex'[sep=0pt -2.5]}}},
                      postaction={decorate}}}
\tikzset{<--/.style = {decoration={markings,
                                  mark=at position 0.001
                                  with {\arrowreversed[scale=2]{latex'[sep=0pt -2.5]}}},
                      postaction={decorate}}}
\tikzset{<->/.style = {decoration={markings,
                                  mark=at position 0.001 with {\arrowreversed[scale=2]{latex'}},
                                  mark=at position 0.999 with {\arrow[scale=2]{latex'}},},
                      postaction={decorate}}}
\tikzset{->-/.style = {decoration={markings,
    mark=at position 0.50+0.225/#1 with {\arrow[scale=2]{latex'}}},
                      postaction={decorate}}}
\tikzset{-<-/.style = {decoration={markings,
                                   mark=at position 0.50-0.225/#1 with {\arrowreversed[scale=2]{latex'}}},
                       postaction={decorate}}}
\tikzset{->>/.style = {decoration={markings,
                                  mark=at position 1 with {\arrow[scale=2]{latex'}}},
                      postaction={decorate}}}
\tikzset{<<-/.style = {decoration={markings,
                                  mark=at position 0 with {\arrowreversed[scale=2]{twolatex'}}},
                      postaction={decorate}}}
\tikzset{<<->>/.style = {decoration={markings,
                                   mark=at position 0 with {\arrowreversed[scale=2]{twolatex'}},
                                   mark=at position 1 with {\arrow[scale=2]{twolatex'}}},
                       postaction={decorate}}}
\tikzset{->>-/.style = {decoration={markings,
                                   mark=at position 0.50+0.450/#1 with {\arrow[scale=2]{twolatex'}}},
                       postaction={decorate}}}
\tikzset{-<<-/.style = {decoration={markings,
                                   mark=at position 0.50-0.450/#1 with {\arrowreversed[scale=2]{twolatex'}}},
                       postaction={decorate}}}

\pgfarrowsdeclare{biggertip}{biggertip}{%
  \setlength{\arrowsize}{1pt}
  \addtolength{\arrowsize}{.1\pgflinewidth}
  \pgfarrowsrightextend{0}
  \pgfarrowsleftextend{-5\arrowsize}
}{%
  \setlength{\arrowsize}{1pt}
  \addtolength{\arrowsize}{.1\pgflinewidth}
  \pgfpathmoveto{\pgfpoint{-5\arrowsize}{4\arrowsize}}
  \pgfpathlineto{\pgfpointorigin}
  \pgfpathlineto{\pgfpoint{-5\arrowsize}{-4\arrowsize}}
  \pgfusepathqstroke
}
\tikzset{
	EdgeStyle/.style = {>=biggertip}
}

\tikzset{circ/.style = {fill, circle, inner sep = 0, minimum size = 3}}
\tikzset{scirc/.style = {fill, circle, inner sep = 0, minimum size = 1.5}}
\tikzset{mstate/.style={circle, draw, black, text=black, minimum width=0.7cm}}

\tikzset{eqpic/.style={baseline={([yshift=-.5ex]current bounding box.center)}}}

\definecolor{mblue}{rgb}{0.2, 0.3, 0.8}
\definecolor{morange}{rgb}{1, 0.5, 0}
\definecolor{mgreen}{rgb}{0, 0.4, 0.2}
\definecolor{mred}{rgb}{0.5, 0, 0}

% topology
\DeclareMathOperator{\dfr}{dfr} % deformation retract
\newcommand{\Cells}{\text{Cells}}
\newcommand{\RP}{\mathbb{RP}}
\newcommand{\CP}{\mathbb{CP}}

% order theory
\newcommand{\join}{\vee}
\newcommand{\meet}{\wedge}

% algebraic geometry
\DeclareMathOperator{\Rad}{Rad} % Radical of modules
\DeclareMathOperator{\Spec}{Spec}
% \renewcommand{\div}{\mathrm{div}} % divisors
\DeclareMathOperator{\Mer}{Mer} % Meromorphic functions
\DeclareMathOperator{\Exc}{Exc} % Exceptional divisor
\DeclareMathOperator{\Td}{Td} % Todd classes

% algebra

\DeclareMathOperator{\trdeg}{tr.deg.}
\DeclareMathOperator{\odd}{odd}
\DeclareMathOperator{\even}{even}
\DeclareMathOperator{\lcm}{lcm}
\DeclareMathOperator{\ord}{ord}
\DeclareMathOperator{\codim}{codim}
\DeclareMathOperator{\Ann}{Ann}
\DeclareMathOperator{\Ext}{Ext}
\DeclareMathOperator{\Sym}{Sym}
\DeclareMathOperator{\Out}{Out}
\DeclareMathOperator{\Aut}{Aut}
\DeclareMathOperator{\End}{End}
\DeclareMathOperator{\Inn}{Inn}
\DeclareMathOperator{\Mat}{Mat}
\DeclareMathOperator{\Fix}{Fix}
\DeclareMathOperator{\std}{std}
\DeclareMathOperator{\sgn}{sgn}
\DeclareMathOperator{\adj}{adj}
\DeclareMathOperator{\id}{id}
\DeclareMathOperator{\op}{op}
\DeclareMathOperator{\tr}{tr}
\DeclareMathOperator{\diag}{diag}
\DeclareMathOperator{\rank}{rank}
\DeclareMathOperator{\rk}{rk}
\DeclareMathOperator{\coker}{coker}
\DeclareMathOperator{\Mult}{Mult} % multiplier algebra
\DeclareMathOperator{\Frac}{Frac} % fraction field
\DeclareMathOperator{\Rat}{Rat}
\DeclareMathOperator{\Gal}{Gal}
\newcommand{\ioplus}{\overset{\text{int}}{\oplus}} % interior direct sum
\newcommand{\GG}{\mathcal G} % Galois correspondence
\newcommand{\FF}{\mathcal F} % Galois correspondence
\newcommand{\cupp}{\smile} % cup product
\newcommand{\capp}{\frown} % cap product
\newcommand{\actson}{\curvearrowright}
\newcommand{\bra}{\left\langle}
\newcommand{\ket}{\right\rangle}
\newcommand{\idealin}{\triangleleft\,}
\newcommand{\normalin}{\triangleleft\,}
\newcommand{\ip}[2]{\left\langle #1, #2 \right\rangle}
\newcommand{\bigslant}[2]
{{\raisebox{.2em}{$#1$}\left/\raisebox{-.2em}{$#2$}\right.}}
\newcommand{\properideal}{%
  \mathrel{\ooalign{$\lneq$\cr\raise.22ex\hbox{$\lhd$}\cr}}}

% categories
\newcommand{\cat}[1]{\mathbf{#1}}
\newcommand{\Ch}{\cat{Ch}} % Chain complexes   
\newcommand{\Set}{\cat{Set}}
\newcommand{\Grp}{\cat{Grp}}
\newcommand{\Ab}{\cat{Ab}}
\newcommand{\Ring}{\cat{Ring}}
\newcommand{\CRing}{\cat{CRing}}
\newcommand{\Top}{\cat{Top}}
\newcommand{\Vect}{\cat{Vect}}
\newcommand{\cmod}{\cat{mod}}

%% matrix groups
\newcommand{\GL}{\mathrm{GL}}
\newcommand{\Or}{\mathrm{O}}
\newcommand{\PGL}{\mathrm{PGL}}
\newcommand{\PSL}{\mathrm{PSL}}
\newcommand{\PSO}{\mathrm{PSO}}
\newcommand{\PSU}{\mathrm{PSU}}
\newcommand{\SL}{\mathrm{SL}}
\newcommand{\msl}{\mathrm{sl}}
\newcommand{\SO}{\mathrm{SO}}
\newcommand{\Spin}{\mathrm{Spin}}
\newcommand{\Sp}{\mathrm{Sp}}
\newcommand{\SU}{\mathrm{SU}}
\newcommand{\U}{\mathrm{U}}
\let\char\relax
\newcommand{\char}{\text{char}}

% analysis
\renewcommand{\Re}{\mathrm{Re}}
\renewcommand{\Im}{\mathrm{Im}}
\NewDocumentCommand{\dx}{o}{
  \IfNoValueTF{#1}
    {\dif x}                  % if no optional argument
    {\dif^{#1} \!x}         % if optional argument
}
\NewDocumentCommand{\dy}{o}{
  \IfNoValueTF{#1}
    {\dif y}                  % if no optional argument
    {\dif^{#1} \!y}         % if optional argument
}
\newcommand{\dt}{\dif t}
\newcommand{\du}{\dif u}
\newcommand{\dv}{\dif v}
\newcommand{\dz}{\dif z}
\newcommand{\ds}{\dif s}
\newcommand{\dw}{\dif w}
\newcommand{\dr}{\dif r}
\newcommand{\dm}{\dif m}
\newcommand{\df}{\dif f}
\newcommand{\dV}{\dif V}
\newcommand{\dS}{\dif S}
\newcommand{\dA}{\dif \mathrm{A}}
\newcommand{\dmu}{\dif \mu}
\newcommand{\dnu}{\dif \nu}
\newcommand{\dtheta}{\dif \theta}
\newcommand{\domega}{\dif \omega}
\newcommand{\dsigma}{\dif \sigma}
\newcommand{\dtau}{\dif \tau}
\newcommand{\dvarphi}{\dif \varphi}
\renewcommand{\dif}{\mathrm{d}}
\renewcommand{\Dif}{\mathrm{D}}
\renewcommand{\triangle}{\Delta}
\newcommand{\ptriangle}{\partial \triangle}
\DeclareMathOperator{\im}{im}
\DeclareMathOperator{\cis}{cis}
\DeclareMathOperator{\rot}{rot}
\DeclareMathOperator{\vspan}{span}
\DeclareMathOperator{\grad}{grad}
\DeclareMathOperator{\curl}{curl}
\DeclareMathOperator{\esssup}{esssup}
\newcommand{\hodge}{{\mathnormal{\star}}}
\DeclareMathOperator{\range}{range}
\DeclareMathOperator{\Hom}{Hom}
\DeclareMathOperator{\Int}{Int}
\DeclareMathOperator{\Conv}{Conv} % convex hull
\newcommand{\ind}{\mathbbm{1}}
\DeclareMathOperator{\Res}{Res}
\DeclareMathOperator{\graph}{graph}
\DeclareMathOperator{\diam}{diam}
\DeclareMathOperator{\supp}{supp}
\DeclareMathOperator{\Vol}{Vol} % Volume
\DeclareMathOperator{\Area}{Area} % Area
\DeclareMathOperator{\area}{area}
\DeclareMathOperator{\Ind}{Ind} % Index
\newcommand{\length}{\mathrm{length}}

% logic
\DeclareMathOperator{\MOD}{MOD}
\DeclareMathOperator{\Theory}{Theory}
\newcommand{\notimplies}{%
  \mathrel{{\ooalign{\hidewidth$\not\phantom{=}$\hidewidth\cr$\implies$}}}}

% nice
\newcommand{\half}{\frac{1}{2}}
\newcommand{\pair}{\del}
\newcommand{\taking}[1]{\xrightarrow{#1}}
\newcommand{\otaking}[1]{\xleftarrow{#1}}
\newcommand{\inv}{^{-1}}
\newcommand{\ot}{\leftarrow}
\newcommand{\ninfty}{-\infty}
\newcommand{\pinfty}{+\infty}
\newcommand{\floor}[1]{\left\lfloor #1 \right\rfloor}
\newcommand{\ceil}[1]{\left\lceil #1 \right\rceil}
\newcommand{\viff}{\Big\Updownarrow} % vertical iff
\newcommand{\bmid}{\,\middle\vert\,} % big mid

% probability
\newcommand{\Prob}{\mathbf{P}}
\renewcommand{\vec}[1]{\boldsymbol{\mathbf{#1}}}
\DeclareMathOperator{\Bin}{Bin}
\DeclareMathOperator{\Geo}{Geo}
\DeclareMathOperator{\Poi}{Poi}
\DeclareMathOperator{\Exp}{Exp}
\DeclareMathOperator{\Var}{Var} % Variance
\DeclareMathOperator{\Cov}{Cov}

% special letters
\newcommand{\N}{\mathbb{N}}
\newcommand{\Z}{\mathbb{Z}}
\newcommand{\Q}{\mathbb{Q}}
\newcommand{\R}{\mathbb{R}}
\newcommand{\C}{\mathbb{C}}
\newcommand{\F}{\mathbb{F}}
\newcommand{\E}{\mathbf{E}}
\newcommand{\D}{\mathbb{D}}
\renewcommand{\H}{\mathbb{H}}
\newcommand{\ps}{\mathcal{P}}
\newcommand{\M}{\mathcal{M}}
\renewcommand{\L}{\mathcal{L}}
\renewcommand{\O}{\mathcal{O}}
\renewcommand{\U}{\mathcal{U}}
\newcommand{\Omicron}{O}
\newcommand{\powerset}{\mathcal{P}}

% text
\newcommand{\st}{\text{ s.t. }}
\newcommand{\tand}{\quad \text{and} \quad}
\newcommand{\tor}{\quad \text{or} \quad}
\newcommand{\stand}{\text{ and }}
\newcommand{\stor}{\text{ or }}
\renewcommand{\tt}[1]{\textnormal{\textbf{(#1).}}} %tt=theorem title GET RID OF

% title format
\title{\textbf{\ncourse}}

\ifx \ntype \notes
  \author{Notes taken by \nauthor\, \nemail \\ Based on lectures by \nlecturer}
  \date{\nterm\ \nyear}
\else
  \ifx \ntype \problems
    \author{These problems were solved by \nauthor \\ \nemail}
  \else
  \fi
\fi

\makeatother

\newcommand{\A}{\mathcal A}

\begin{document}
\maketitle

% Insert cool image here

\newpage
\tableofcontents
\newpage

\section{Introduction}

\subsection{Motivation}
The Riemann integral we have known so far is fairly limited.
For example it doesn't allow us to compute the Riemann integral
of Dirichlet's function $f \colon [0,1] \to \R$ defined as
\[
  f(x) = \mathbbm{1}_{\Q \cap [0,1]} =
  \begin{cases}
    1, &x \in \Q \cap [0,1] \\
    0, &\text{otherwise}
  \end{cases}.
\]
In his thesis Lebesgue introduced a new type of integral called
a Lebesgue integral that allows us to compute integrals for functions
like Dirichelt's function, and he continued to develop more concepts
like measure, and almost everywhere.

\subsection{Motivation for Lebesgue integral}
Let $f(x) = \mathbbm{1}_A$ be the function that 
we want $\int \mathbbm{1}_A$ to be the volume of the set $A$.

First we would like we define what is a volume of a set.
We would want to require a couple of things
\begin{enumerate}
  \item[(1)] $\mu(A)$ is defined for all $A \subseteq \R^n$;
  \item[(2)] $\mu([0,1]^n) = 1^n = 1$;
  \item[(3)] $\mu$ to be invariant to congruations (isometries).
  \item[(4)] If $\set{A_i}_{i=1}^{\infty}$ is a countable sequence
    of pairwise disjoint sets then
    \[
      \mu\del{\bigcup_{i=1}^{\infty} A_i} = \sum_{i=1}^{\infty} \mu(A_i).
    \]
\end{enumerate}

\begin{remark}
  Property $(4)$ is called $\sigma$-additivity.
\end{remark}

\begin{theorem}[Hausdorrf, 1914]
  There is no function that satisfies $(1) - (4)$ at the same time.
\end{theorem}

We will prove this theorem later.
For now we can only try to weaken the requirements.
For example instead of $\sigma$-additivity we might require finite
additivity.

\begin{theorem}
  There exists a function that satisfies the wanted requirements in
  dimenstions $1$ and $2$ but not in dimension $n \geq 3$.
\end{theorem}

\begin{remark}
  For example in $n = 3$ we have the Banach--Tarski paradox.
\end{remark}

\begin{paradox}[Banach--Tarski, 1924]
  For every $n \geq 1$ we can divide $S^2$ in $\R^n$ to a finite amount 
  of parts such that when they are rotated and rearranged,
  can form a new sphere of any desired size.
\end{paradox}
\begin{remark}
  The bigger sphere we would want to form, the greater is the minimal
  pieces we need to divide the unit sphere in order to form it.
\end{remark}

This paradox is based on the use of the axiom of choice, but since we
assume the axiom of choice in this course we still need to modify the
requirements.

Instead of $(4)$ we require subadditivity:
\[
\mu\del{\bigcup_{i=1}^{\infty} A_i} \le \sum_{i=1}^{\infty} \mu(A_i).
\]
This is possible but instead we would like to keep $\sigma$-additivity
but instead give up on requirement $(0)$.
We will only define volume only for ``nice'' sets which include
all of the sets we work with in analysis or geometry etc.\ these
sets form a $\sigma$-algebra on $\R^n$.

\newpage

\section{Algebras and \texorpdfstring{$\sigma$}{s}-algebras}

\subsection{Definitions}

\begin{definition}[Algebra]
  An algebra on a nonempty set $X$ is a collection $\mathcal A$ of
  subsets of $X$ such that
  \begin{enumerate}
    \item[(1)] $\emptyset, X \in \mathcal A$;
    \item[(2)] $\mathcal A$ is closed under complements;
    \item[(3)] $\mathcal A$ is closed under finite unions and intersections.
  \end{enumerate}
\end{definition}

\begin{definition}[$\sigma$-algebra]
  A $\sigma$-algebra on a nonempty set $X$ is an algebra $\mathcal A$ on
  $X$ that is also closed under countable unions and intersections.
\end{definition}

\begin{remark}
  A $\sigma$-algebra is sometimes called a $\sigma$-field.
  That is why it is sometimes denoted $\mathcal F$.
\end{remark}

\begin{remark}
  From De-Morgan laws we know that
  \[
    \del{\bigcup_{i \in I} A_i}^c =
    \bigcap_{i \in I} A_i^c
  \]
  so it is only necessary to require closure under countable unions
  or countable intersections.
\end{remark}

\begin{remark}
  It is also possible to require $\mathcal A$ to be nonempty instead of 
  $(1)$.
  Then for $A \in \mathcal A$ we have
  \[
    A \in \A \implies A^c \in \A \implies
    A \cup A^c = X \in \A \stand
    A \cap A^c = \emptyset \in \A.
  \]
\end{remark}

\begin{example}
  $\A = \set{\emptyset, X}$ and $\A = 2^X$ are the smallest
  and biggest $\sigma$-algebras on $X$ respectively.
\end{example}

\begin{example}
  If $X$ is not countable.
  Then
  \[
    \mathcal A =
    \set{E \subseteq X \colon E \stor E^c \text{ are countable}} \neq
    2^X
  \]
  is a $\sigma$-algebra.
\end{example}

\begin{definition}[Cocountablility]
  Let $X$ be a set.\ Then $A$ is called cocountable if $A^c$ is countable.
\end{definition}

\begin{example}[Generated $\sigma$-algebra]
  Let $F \subseteq 2^X$ be a family of subsets of $X$.
  The $\sigma$-algebra generated by $F$ is defined as
  \[
    \sigma(F) = \bigcap \set{\mathcal A \subseteq 2^X \colon 
    \A \text{ is a $\sigma$-algebra} \stand F \subset \A}.
  \]
\end{example}
\begin{remark}
  Notice that the intersection is indeed a $\sigma$-algerba.
\end{remark}
\begin{remark}
  Let $\A$ be a $\sigma$-algebra and $F \subseteq \A$.
  Then $\sigma(F) \subseteq A$.
\end{remark}
\begin{corollary}
  Suppose $F_1 \subseteq \sigma(F_2)$ and $F_2 \subseteq \sigma(F_1)$.
  Then $\sigma(F_1) = \sigma(F_2)$.
\end{corollary}

\begin{definition}[Borel $\sigma$-algebra]
  Let $X$ be a topological space.
  Then we define the Borel $\sigma$-algebra as the $\sigma$-algebra generated
  by the open sets in $X$.
  \[
    B(X) = \sigma\del{\set{G \subset X \mid G \text{ is open}}}.
  \]
\end{definition}
\begin{remark}
  Recall that $G$ denotes an open set, $F$ a closed set, $G_{\delta}$
  a countable intersection of open sets and $F_{\sigma}$ a countable
  union of closed sets.
  Similarly $G_{\delta \sigma}$ is a countable union of $G_{\delta}$ sets
  etc.
\end{remark}
\begin{remark}
  \[
    B(X) = \sigma\del{\set{F \subset X \mid F \text{ is closed}}}.
  \]
\end{remark}

\begin{proposition}
  $B(\R)$ is generated by the collection of any type of interval.
\end{proposition}
\begin{proof}
  Consider the collection of open intervals in $\R$.
  Since any open set in $\R$ is a countable union of disjoin open intervals,
  we get that $B(\R) = \sigma\del{(a,b)}$.

  Consider the collection of closed intervals in $\R$.
  Since any closed interval is a countable intersection of open sets,
  we have that $[a,b] \subseteq \sigma\del{(a,b)}$.
  Since any open interval is a countable union of closed sets,
  we have that $(a,b) \subseteq \sigma\del{[a,b]}$.
  % This shows that $\sigma\del{[a,b]} = \sigma\del{(a,b)} = $B(\R)$
  
  The proof for other types of intervals (like $(a,b]$) is similar and thus
  omitted.
\end{proof}

\subsection{The product \texorpdfstring{$\sigma$}{s}-algebra}

\begin{definition}[Product $\sigma$-algebra]
  Let $\set{\A_{i}}_{i}$ be a collection of $\sigma$-algebras on
  $\set{X_{i}}_{i}$.
  Then, the product $\sigma$-algebra $\bigotimes_{i \in I} \A_{i}$ is
  the $\sigma$-algebra generated by the cylindrical sets:
  \[
    \mathcal{S} = \set{\prod_{i \in I}{U_{\alpha}} \mid 
    \exists j \in I \st 
    \text{ $U_{i} = X_{i}$ 
    for $i \in I \setminus \{j\}$ and $U_j \in \A_j$}},
  \]
  which is the also the set that generates the product topology on
  the product space.
\end{definition}
\begin{remark}
  In the case of $|I| < \aleph_0$ this is equivalent to the $\sigma$-algebra
  generated by the open sets in $\bigotimes_{i \in I} \A_i$ with respect to the
  box topology.
\end{remark}

\begin{proposition}
  $B(\R^n) = \bigotimes_{i=1}^{n} B(\R)$.
\end{proposition}
\begin{proof}
  Showing that $\prod_{i=1}^{n} (a_i, b_i) \subseteq B(\R^n)$ can be
  done simply by definition.
  We can also show that $B(\R^n) \subseteq \prod_{i=1}^{n} (a_i, b_i)$
  because for an open set $A \subseteq \R^n$ we
  can see that it is the countable union of all the open boxes with
  rational sized edges contained in $A$ around any rational point $q \in A$
  (which exist because $A$ is open).
\end{proof}

\newpage

\section{Measure}
\subsection{Definitions}

\begin{definition}[Measure]
  Let $\A$ be a $\sigma$-algebra on $X$.
  A function $\mu \colon \A \to [0,\infty]$ is called a measure if
  \begin{enumerate}
  \item[(1)] $\mu(\emptyset) = 0$;
  \item[(2)] Given a sequence $\set{E_i}_{i=1}^{\infty} \subseteq \A$ of
    disjoint sets we have
    \[
      \mu\del{\bigsqcup_{i=1}^{\infty} E_i} =
      \sum_{i=1}^{\infty} \mu(E_i).
    \]
  \end{enumerate}
\end{definition}
\begin{remark}
  Property $(2)$ is called $\sigma$-additivity.
  It obviouly implies finite additivity because we can choose 
  $E_j = \emptyset$.
\end{remark}

\begin{definition}[Measurable space]
  The space $(X,\A)$ is called a measurable space.
  The elements of $\A$ are called measurable sets.
\end{definition}

\begin{definition}[Measure space]
  The triple $(X,\A,\mu)$ is called a measure space.
\end{definition}

\begin{definition}[Finite measure]
  A finite measure is a measure $\mu$ on $X$ such that $\mu(X) < \infty$.
\end{definition}

\begin{definition}[$\sigma$-finite measure]
  A $\sigma$-finite measure is a measure $\mu$ on $X$ such that
  $X = \bigcup_{i=1}^{\infty} E_i$ for $E_i \in \A$ such that
  $\mu(E_i) < \infty$.
\end{definition}

\begin{definition}[Borel measure]
  A measure $\mu$ on a topological space $X$ is called a Borel measure
  if $B(X) \subseteq \A$.
\end{definition}

\begin{example}[Delta measure]
  Let $x_0 \in X$ and $\A = 2^X$.
  Then the delta measure is
  \[
    \delta_{x_0}(E) =
    \begin{cases}
      1, &x_0 \in E \\
      0, &x_0 \notin E
    \end{cases}.
  \]
\end{example}

\begin{example}[Counting measure]
  The counting measure is the measure on $\A = 2^X$ such that
  $\mu(E) = |E|$.
\end{example}

\begin{example}
  \label{exmp:relevency-measure-on-r}
  Suppose $\aleph_0 < |X|$ and let $\A$ be the $\sigma$-algebra of the countable
  or cocountable subsets of $X$. Then the follwing is a measure on $X$,
  \[
    \mu(E) =
    \begin{cases}
      0, &|E| \le \aleph_0 \\
      1, &|E^c| = \aleph_0
    \end{cases}.
  \]
\end{example}

\begin{example}
  Suppose $|X| = \infty$ and $\A = 2^X$.
  Then the following is a measure on $X$,
  \[
    \mu(E) =
    \begin{cases}
      0, &|E| < \infty \\
      1, &|E| = \infty
    \end{cases}.
  \]
\end{example}
\begin{remark}
  The last example is a finite-additive, but not $\sigma$-additive.
\end{remark}

\begin{proposition}
  Let $(X, \A, \mu)$ be a measure space.
  Then
  \begin{enumerate}
    \item[(1)] Let $E,F \in \A$ such that $E \subseteq F$.
      Then $\mu(E) \le \mu(F)$.
    \item[(2)] Let $\set{E_i}_{i=1}^{\infty} \subseteq A$.
      Then
      \[
        \mu\del{\bigcup_{i=1}^{\infty} E_i} \le \sum_{i=1}^{\infty} \mu(E_i)
      \]
    \item[(3)] Let $\set{E_i}_{i=1}^{\infty} \subseteq A$ such that
      $E_1 \subseteq E_2 \subseteq \dots$ then
      \[
        \mu\del{\bigcup_{i=1}^{\infty} E_i} =
        \lim_{i \to \infty} \mu(E_i)
      \]
    \item[(4)] Let $\set{E_i}_{i=1}^{\infty} \subseteq A$ such that
      $E_1 \supseteq E_2 \supseteq \dots$ then
      \[
        \mu\del{\bigcap_{i=1}^{\infty} E_i} =
        \lim_{i \to \infty} \mu(E_i)
      \]
  \end{enumerate}
\end{proposition}
\begin{proof} \phantom{}
\begin{enumerate}
  \item[(1)] We have that $F = E \sqcup (F \setminus E)$ so
    \[
      \mu(F) = \mu(E) + \mu(F \setminus E) \geq \mu(E).
    \]
  \item[(2)] Define the sets
    \begin{align*}
      F_1 &:= E_1 \\
      F_2 &:= E_2 \setminus E_1 \\
      &\vdotswithin{:=} \\
      F_k &:= E_k \setminus \del{\bigcup_{j=1}^{k-1} E_j}.
    \end{align*}
    We now have that $\set{F_k}_{k=1}^{\infty}$ are disjoint and
    for all $n \geq 1$ we have
    \[
      \bigsqcup_{j=1}^{n} F_j =
      \bigcup_{j=1}^{n} E_j
    \]
    which implies that
    \[
      \mu\del{\bigcup_{j=1}^{\infty} E_j} =
      \mu\del{\bigsqcup_{j=1}^{\infty} F_j} =
      \sum_{j=1}^{\infty} \mu(F_j) \le
      \sum_{j=1}^{\infty} \mu(E_j)
    \]
    where the last inequality follows from $F_j \subseteq E_j$ and $(1)$.
  \item[(3)] Define similarly
    \begin{align*}
      F_0 &:= \emptyset \\
      F_1 &:= F_1 \setminus E_0 \\
      &\vdotswithin{:=} \\
      F_k &:= F_k \setminus F_{k-1}.
    \end{align*}
    We now have that $\set{F_k}_{k=1}^{\infty}$ are disjoint and
    for all $n \geq 1$ we have
    \[
      \bigsqcup_{j=1}^{n} F_j =
      \bigcup_{j=1}^{n} E_j
    \]
    which implies that
    \[
      \mu\del{\bigcup_{j=1}^{n} E_j} =
      \mu\del{\bigsqcup_{j=1}^{n} F_j} =
      \sum_{j=1}^{\infty} \mu(F_j) =
      \lim_{n \to \infty} \sum_{j=1}^{n} \mu(F_j) =
      \lim_{n \to \infty}
      \underbrace{\mu\del{\bigsqcup_{j=1}^{n} F_j}}_{\mu(E_n)}
    \]
    where the last euqality follows from additivity.
  \item[(4)] This can be concluded using De Morgan's laws on $(3)$.
\end{enumerate}
\end{proof}

\begin{definition}[$\mu$-negligible set]
  A negligible set is a set $E \in \A$ such that $\mu(E) = 0$.
\end{definition}

\begin{definition}[Almost everywhere]
  Let $(X,\A,\mu)$ be a measure space.
  We say that a certain property is true \textit{almost everywhere} if it is
  true for any $x \in X \setminus E$ such that $E$ is of measure zero
  ($\mu$-negligible).
\end{definition}

\begin{example}
  We have that $\sin x \neq 0$ almost everywhere because
  \[
    |E| = \left|\set{x \in \R \mid \sin x = 0}\right| = \aleph_0
  \]
  is $\mu$-negligible for the measure $\mu$ from
  \Cref{exmp:relevency-measure-on-r} with $X = \R$ defined
  \[
    \mu(E) =
    \begin{cases}
      0, &|E| \le \aleph_0 \\
      1, &|E^c| = \aleph_0
    \end{cases}.
  \]
\end{example}
\begin{remark}
  A countable union of $\mu$-negligible sets is $\mu$-negligible from
  subadditivity.
\end{remark}
\begin{remark}
  A subset of a $\mu$-negligible set is $\mu$-negligible since measures are
  monotone.
\end{remark}

\subsection{Complete measure spaces}

\begin{definition}[Complete measure space]
  A triple $(X,\A,\mu)$ is called a complete measure space if
  for any $F \in \A$, if $\mu(F) = 0$ then for all $E \subseteq F$
  we have $E \in \A$.
\end{definition}

\begin{theorem}[Completion theorem]
  Let $(X, \A, \mu)$ be a measure space.
  Define
  \begin{align*}
    \mathcal N &:= \set{N \in \A \mid \mu(N) = 0} \\
    \overline \A &:= 
    \set{E \cup F \mid E \in \A \stand F \subseteq N \in \mathcal N}.
  \end{align*}
  Then $\overline A$ is a $\sigma$-algebra, and there exists 
  a unique extension $\bar \mu$ of $\mu$ such that $(X, \overline \A, \bar \mu)$
  is a complete measure space.
\end{theorem}
\begin{proof}
  In the problems section.
\end{proof}
\begin{remark}
  $\overline A$ is the smallest $\sigma$-algebra containing $\A$ and all
  the subsets of $\mu$-negligible sets.
\end{remark}

\noindent
Our goal is to construct a measure space $(\R^n, \A, \mu)$ such that
\[
  \mu\del{\prod_{i=1}^{n} [a_i,b_i]} =
  \prod_{i=1}^{n} |a_i - b_i|.
\]
Recall that earlier we saw that
\[
  B(\R^n) =
  \bigotimes_{i=1}^{n} B(\R) =
  \sigma\del{\prod_{i=1}^{n} [a_i,b_i]}
\]
which implies that $B(\R^n) \subseteq \A$.
To define the measure for other sets we can use an outer measure
\[
  \mu^*(E) := \inf \set{\sum_{i=1}^{\infty} \mu(R_i) \mid E \subseteq 
  \bigcup_{i=1}^{n} R_i}.
\]
We now have two questions.
\begin{itemize}
  \item[(1)] Do we always have $\mu^*(R) = \mu(R)$? Yes.
  \item[(2)] Is $\mu^*$ a measure? No, it is an outer measure.
\end{itemize}
\begin{remark}
  Defining $\mu^*$ only on ``good'' sets gives us a measure called
  Lebesgue measure.
  Good sets are sets for which the outer measure is equal to the inner measure.
\end{remark}

\begin{definition}[Outer measure]
  Let $X \neq \emptyset$. Then $\mu^* \colon 2^X \to [0,\infty]$ is
  called an outer measure if
  \begin{enumerate}
    \item[(1)] $\mu^*(\emptyset) = 0$;
    \item[(2)] $\mu^*$ is monotone (if $A \subset B$ then $\mu^*(A) < \mu^*(B)$);
    \item[(3)] $\mu^*$ is subadditive
      $\del{\mu^*\del{\bigcup_{i=1}^{\infty} E_i} \le 
      \sum_{i=1}^{\infty} \mu^*(E_i)}$.
  \end{enumerate}
\end{definition}

\begin{proposition}
  Let $\mathcal E \subseteq 2^X$ be a family of subsets of $X$ and
  $\varphi \colon \mathcal E \to [0,\infty]$ such that
  $\emptyset, X \in \mathcal E$ and $\varphi(\emptyset) = 0$.
  For all $A \subseteq X$ we define
  \[
    \mu^*(A) = \inf\set{\sum_{i=1}^{\infty} \varphi(E_i) \mid 
    E_i \in \mathcal E \stand A \subseteq \bigcup_{i=1}^{\infty} E_i}
  \]
  then $\mu^*$ is an outer measure.
\end{proposition}
\begin{remark}
  It is not promised that $\mu^*(E) = \varphi(E)$ for $E \in \mathcal E$.
\end{remark}
% ADD EXAMPLE?
\begin{proof} \phantom{}
\begin{enumerate}
  \item[(1)] We can cover $\emptyset$ by $E_i = \emptyset$.
  \item[(2)] Let $A \subseteq B$.
    We can use the cover of $B$ for $A$ since 
    $A \subseteq B \subseteq \bigcup E_i$ which implies that
    $\mu^*(A) \le \mu^*(B)$.
  \item[(3)] By the definition of the infimum there exists a covering
    $A_i \subseteq \cup_{k=1}^{\infty} E_k^i$ such that
    \[
      \mu\del{\bigcup_{k=1}^{\infty} E_k^i} \le 
      \sum_{i=1}^{\infty} \mu^*(A_i) + \epsilon \cdot 2^{-i}
    \]
    which implies that 
    $\bigcup_{i=1}^{\infty} A_i \subseteq \bigcup_{k,i=1}^{\infty} E_k^i$
    and so
    \[
      \mu^*\del{\bigcup_{i=1}^{\infty} E_i} \le 
      \sum_{k,i}^{\infty} \varphi(E_k^i) \le
      \sum_{i=1}^{\infty} \mu^*(E_i) + \epsilon
    \]
    but $\epsilon > 0$ is arbitrary which completes the proof.
\end{enumerate}
\end{proof}

\begin{remark}
  Every outer measure can be constructed by $\varphi$ in this way.
\end{remark}

\begin{definition}[$\mu^*$-measurable set]
  A set $A \subseteq X$ is called $\mu^*$-measurable if for all $E \subseteq X$,
  \[
    \mu^*(E) = \mu^*(E \cap A) + \mu^*(E \cap A^c).
  \]
\end{definition}

\begin{remark}
  It is clear that the $\le$ direction is always satisfied from subadditivity,
  so it is only necessary to verify the $\geq$ direction.
  In particular, it is only necessary to verify it for sets such that
  $\mu^*(E) < \infty$.
  Let $A \subseteq E$ be such that $\mu^*(E) < \infty$.
  Then we need to verify that
  \[
    \underbrace{\mu^*(A)}_{\text{outer measure}} = 
    \underbrace{\mu^*(E) - \mu^*(E \cap A^c)}_{\text{inner measure}}.
  \]
\end{remark}

\subsection{Carath\'eodory's outer measure theorem}

\begin{theorem}[Carath\'eodory's outer measure theorem]
  \label{thm:car-outer-measure}
  Let $\mu^*$ be an outer measure on $X$. Then
  \begin{enumerate}
    \item[(1)] $\mathcal F_{\mu^*} := 
      \set{A \subseteq X \mid A \text{ is }\mu^*\text{-measurable}}$ is a 
      $\sigma$-algebra.
    \item[(2)] The triplet
      $(X, \mathcal F_{\mu^*}, \mu^*\vert_{\mathcal F_{\mu^*}})$
      is a complete measure space.
  \end{enumerate}
\end{theorem}
\begin{remark} \label{recitations-one}
  We saw in the recitations that in order
  to show that $\mathcal F$ is a $\sigma$-algebra, it suffices to prove that
  if $E_i \in \mathcal F$ are disjoint, then
  \[
    \bigsqcup_{i=1}^{\infty} E_i \in \mathcal F.
  \]
\end{remark}
\begin{proof} \phantom{}
\begin{enumerate}
  \item[(1)] Let us show that $\mathcal F = \mathcal F_{\mu^*}$ is an algebra.
  $\mathcal F$ is closed under the complement operation because by the
  symmetric definition, if $A$ is $\mu^*$-measurable so is $A^c$.

  Let $A,B \in \mathcal F$ and $E \subseteq X$.
  Since $A$, $B$ are $\mu^*$-measurable we have that
  \begin{align*}
    \mu^*(E) &= \mu^*(E \cap A) + \mu^*(E \cap A^c) \\ &=
    \mu^*(E \cap A \cap B) + \mu^*(E \cap A \cap B^c) +
    \mu^*(E \cap A^c \cap B) + \mu^*(E \cap A^c \cap B^c).
  \end{align*}
  We also have that
  \[
    (A \cup B) \cap E =
    \del{(A \cap B) \cap E} \sqcup
    \del{(A^c \cap B) \cap E} \sqcup
    \del{(A \cap B^c) \cap E}.
  \]
  Thus, from subadditivity
  \[
    \mu^*(E) \geq \mu^*(E \cap (A \cup B)) \mu^*(E \cap (A \cup B)^c)
  \]
  which implies $A \cup B \in \mathcal F$.
  By induction we get that  $\mathcal F$ is closed under finite unions
  which makes it an algebra.
  Moreover, let $A,B \in \mathcal F$ be disjoint sets. Then
  \[
    \mu^*(\underbrace{A \sqcup B}_{E}) =
    \mu^*((A \sqcup B) \cap A) +
    \mu^*((A \sqcup B) \cap A^c) =
    \mu^*(A) + \mu^*(B)
  \]
  so $\mu$ is a finite additive measure on $\mathcal F$.
  
  Let $\set{A_j}_{j=1}^{\infty} \subseteq \mathcal F$ be a sequence of
  disjoint sets.
  Define
  \[
    B_n := \bigsqcup_{i=1}^{n} A_i \in \mathcal F
    \tand
    B_{\infty} := \bigsqcup_{i=1}^{\infty} A_i.
  \]
  For all $E \subseteq X$ we have
  \begin{align*}
    \mu^*(E \cap B_n) &= 
    \mu^*((E \sqcup B_n) \cap A_n) +
    \mu^*((E \sqcup B_n) \cap A_n^c) \\ &=
    \mu^*(E \cap A_n) + \mu^*(E \cap B_{n-1}) \\ &=
    \sum_{i=1}^{n} \mu^*(E \cap A_i).
  \end{align*}
  Using this and since $\mu^*$ is monotone we have
  \[
    \mu^*(E) =
    \mu^*(E \cap B_n) + \mu^*(E \cap B_n^c) \geq
    \sum_{i=1}^{n} \mu^*(E \cap A_i) + \mu^*(E \cap B_{\infty}^c).
  \]
  Taking $n \to \infty$ we have
  \[
    \mu^*(E) \geq
    \sum_{i=1}^{\infty} \mu^*(E \cap A_i) + \mu^*(E \cap B_{\infty}^c) \geq
    \mu^*\del[4]{E \cap \underbrace{\bigcup_{i=1}^{\infty} A_i}_{B_{\infty}}}   
    + \mu^*(E \cap B_{\infty}^c) \geq
    \mu^*(E)
  \]
  which implies that $B_{\infty} \in \mathcal F$ ($\sigma$-additive).
  From \Cref{recitations-one} we have that $\mathcal F$ is a $\sigma$-algebra.
  \item[(2)] Let $A \subseteq X$ such that $\mu^*(A) = 0$.
  Then we have
  \[
    \mu^*(E) \le
    \mu^*(E \cap A) + \mu^*(E \cap A^c) \le
    \underbrace{\mu^*(A)}_{0} + \mu^*(E) =
    \mu^{E}
  \]
  which implies that all the expressions must be equal so
  \[
    \mu^*(E) =
    \mu^*(E \cap A) + \mu^*(E \cap A^c)
  \]
  which means $A \in \mathcal F_{\mu^*}$.
  Since $\mathcal F_{\mu^*}$ contains all $\mu^*$-negligible sets, it is
  complete which also completes the proof.
\end{enumerate}
\end{proof}

% Applications of caratheodory outer measure theorem.
\subsubsection{Constructing Lebesgue measure on \texorpdfstring{$\R^n$}{Rn}}

Let
\[
  \mathcal R =
  \set{R = \prod_{i=1}^{\infty} [a_i,b_i] \middle\vert a_i,b_i \in \R},
\]
and define
\[
  \varphi(R) = \Vol(R) = \prod_{i=1}^{\infty} |a_i - b_i|
  \tand
  \varphi(\R) = \infty.
\]
Let us define the outer Lebesgue measure for $E \subseteq \R^n$
\[
  m^*(E) = \inf\set{\sum_{i=1} \Vol(R_i) \middle\vert
  E \subseteq \bigcup_{i=1}^{\infty} R_i \st R_i \in \mathcal R}.
\]
We saw that it is an outer measure.
By \Cref{thm:car-outer-measure} the space 
$(\R^n,  \mathcal F_{m^*}, m^*\vert_{\L(\R^n)})$ is a complete measure space.
We say that
\[
  m^*\vert_{\L(\R^n)} = m \quad 
  \parbox{8em}{ \phantom{calle}is called \\ Lebesgue measure}
  \tand
  \mathcal F_{m^*} = \L(\R^n) \quad
  \parbox{10em}{is the $\sigma$-algebra \\ of the Lebesure \\
  measurable sets in $\R^n$.}
\]
We need to show that
\begin{enumerate}
  \item[(1)] For all $R \in \mathcal R$ we have $m^*(R) = \Vol(R)$;
  \item[(2)] $\mathcal R \subseteq \L(\R^n) \iff 
    B(\R^n) = \sigma(\mathcal R) \subseteq \L(\R^n)$.
\end{enumerate}
\begin{remark}
  This means that every open set (in particular every Borel set)
  is Lebesgue measurable, and in particular that $m(R) = m^*(R) = \Vol(R)$.
\end{remark}

\begin{proof}
  $(1), (2)$ To be added.
\end{proof}

\subsubsection{Construction of Lebesgue--Stieltjes measure}
Given a nondecreasing monotonic function $F \colon \R \to \R$ which is
continuous from the right, we would like to construct a Borel measure $\mu_F$
on $\R$, such that $\mu_F((a,b]) = F(b) - F(a)$.

The motivation behind this construction is that given a finite
Borel measure $\mu$ we can construct $F(x) = \mu\del{(\ninfty,x]}$
we can easily see that
\[
  \mu\del{(a,b]} =
  \mu\del{(\ninfty,b] \setminus (\ninfty,a]} = F(b) - F(a).
\]
The function $F$ is clearly monotonic and continuous from the right.
We have shown the direction $\mu \to F$ which is the motivation for
the direction $F \to \mu$.

First we define an outer measure $\mu_{F}^{*} \colon 2^{\R} \to [0,\infty]$
where $\mathcal E = \set{(a,b] \mid a < b)}$ and
using the function $\varphi\del{(a,b])} = F(b) - F(a)$ so that
\[
  \mu_{F}^{*}\del{(E])} =
  \inf\set{\sum_{i=1}^{\infty} \varphi\del{(a,b])}
    \middle\vert E \subseteq \bigcup_{i=1}^{\infty} (a_i,b_i]} =
  \inf\set{\sum_{i=1}^{\infty} \del{F(b_i) - F(a_i)}
    \middle\vert E \subseteq \bigcup_{i=1}^{\infty} (a_i,b_i]}.
\]
According to \Cref{thm:car-outer-measure} the triplet
$(\R, \mathcal L_{F}, \mu_{F}^{*}\vert_{\mathcal L_{F}})$
is a complete measure space, where $\mathcal L_{F}$ is the $\mu_{F}^{*}$
measurable sets.
We denote $\mu_{F} := \mu_{F}^{*}\vert_{\mathcal L_{F}}$.
We still need to show that
\begin{enumerate}
  \item[(1)] for all $a < b$ that
    $\mu_{F}\del{(a,b])} = \mu^{*}\del{(a,b]} = F(b) - F(a)$;
  \item[(2)] $(a,b] \in \mathcal L_{F} \iff \mathcal B(\R) = \sigma\del{(a,b]} \subseteq L_F \iff (\ninfty,b) \in \mathcal L_{F}$.
\end{enumerate}

\subsection{Pre-measure}

\begin{definition}[Pre-measure]
  Let $\A_0$ be an algebra on $X$.
  A function $\mu_0 \colon \A_0 \to [0,\infty]$ is called a pre-measure if
  \begin{enumerate}
    \item[(1)] $\mu_0(\emptyset) = 0$.
    \item[(2)] If $\set{E_i}_{i=1}^{\infty} \subseteq \A_0$ are disjoint
      and $\bigsqcup_{i=1}^{\infty} E_i \in \A_0$, then
      \[
        \mu_0\del{\bigsqcup_{i=1}^{\infty} E_i} =
        \sum_{i=1}^{\infty} \mu_0(E_i).
      \]
  \end{enumerate}
\end{definition}
\begin{remark}
  A pre-measure is not a measure because is not necessarily defined on
  a $\sigma$-algebra.
\end{remark}
\begin{remark}
  In particular we get that $\mu_0$ is finite-additive and monotone.
\end{remark}
\begin{remark}
  If $\mu$ is a measure on a $\sigma$-algebra $\A$,
  then $\mu\vert_{\A_0}$ is a pre-measure for every subalgebra
  $\A_0 \subset \A$.
\end{remark}

\begin{theorem}[Carath\'eodory's extension theorem]
  \label{thm:car-extension-theorem}
  Let $\mu_0$ be a pre-measure on an algebra $\A_0$,
  and let $\A = \sigma(\A_0)$.
  Then there exists a measure $\mu$ on $\A$ such that $\mu\vert_{\A_0} = \mu_0$
  defined by $\mu = \mu^*\vert_{\A}$ where
  \[
    \mu^*(E) = \inf \set{\sum_{i=1}^{\infty} \mu_0(R_i) \mid E \subseteq 
    \bigcup_{i=1}^{n} E_i \stand E_i \in \A_0}
  \]
  is the outer measure induced by $\mu_0$.
  Moreover, if $\mu_0$ is $\sigma$-finite, that is
  $X = \bigcup_{j=1}^{\infty} X_j$, $X_j \in \A_0$, and $\mu_0(X_j) < \infty$,
  then the extension is unique (and the existence and uniquness are
  for $F_{\mu^*}$).
\end{theorem}

% "In other words"

\begin{corollary}
  If $\mu_1$ and $\mu_2$ are measures on $(X,\A)$, $\A = \sigma(\A_0)$ for
  some algebra $\A_0 \subset \A$,
  $\mu_1\vert_{\A_0} = \mu_2\vert_{\A_0}$, and $\mu_1$, $\mu_2$ are both
  $\sigma$-finite, then $\sigma_1 = \sigma_2$.
\end{corollary}
\begin{proof}
  To be added
\end{proof}

\begin{corollary}
  Every Borel measure $\mu$ on $\R$ that is ``locally finite''
  (for all $M$ we have $\mu([-M,M]) < \infty$) is of the form $\mu = \mu_F$
  where $\mu_F$ is Stieltjes' measure for some nondecreasing function
  $F \colon \R \to \R$ that is continuous from the right.
  If $\mu$ is finite, then $F(x) = \mu((-\infty,x])$,
  more generally:
  \[
    F(x) =
    \begin{cases}
      \mu((0,x]), &x > 0 \\
      0, &x = 0 \\
      -\mu((x,0]), &x < 0
    \end{cases}.
  \]
\end{corollary}
\begin{proof}
  To be added
\end{proof}

\begin{example}
\end{example}

We can now go back to prove \Cref{thm:car-extension-theorem}.
\begin{proof}
\end{proof}

% Some remarks missing

\begin{remark}
  We have that $m\del{\prod_{i=1}^{n}(a_i,b_i)} = \prod_{i=1}^{n}(a_i,b_i)$
  because we can approximate each edge with $[a_i + 1/k, b_i - 1/k]$ from
  inside.
  We also have for boxes $\set{R_i}_{i=1}^{\infty} \subset \mathcal R$
  \[
    m\del{\bigsqcup_{i=1}^{\infty} R_i} =
    \sum_{i=1}^{\infty} \mu(R_i).
  \]
\end{remark}

\begin{remark}
  For every $x \in \R^n$ we have $m(\set{x}) = 0$ by definition.
  This also implies that $m(\Q^n) = 0$ from $\sigma$-additivity.
\end{remark}

\begin{remark}
  Recall again that $m(A) = m^*(A) = 0$ for $A \subseteq \R^n$
  if and only if
  \[
    \forall \epsilon \quad \exists \set{R_i}_{i=1}^{\infty} \subseteq \mathcal R
    \st
    A \subseteq \bigcup_{i=1}^{\infty} R_i \stand
    \sum_{i=1}^{\infty} \Vol(R_i) < \epsilon.
  \]
\end{remark}

\subsection{The relation between \texorpdfstring{$\L(\R^n$)}{L} and
  \texorpdfstring{$B(\R^n)$}{B}}
\begin{definition}[Regular measure]
  A Borel measure on $(\R^n, \A)$ for any metric or topological space
  is said to be
  \begin{itemize}
    \item Inner regular if
      \[
        \forall E \in \A \quad
        \mu(E) = \sup\set{\mu(K) \mid
          K \subseteq E \stand \text{$K$ is compact and measurable}}.
      \]
    \item Outer regular if
      \[
        \forall E \in \A \quad
        \mu(E) = \inf\set{\mu(G) \mid
          E \subseteq G \stand \text{$G$ is open and measurable}}.
      \]
  \end{itemize}
\end{definition}

\begin{proposition}
  \label{prp:implies-regular}
  Given a Borel measure $\mu$ on $\R^n$ of the form
  $\mu = \mu^*\vert_{\A_{\mu^*}}$ for
  \[
    \mu^*(E) = \inf\set{\sum_{i=1}^{\infty} \mu(G_i) \mid E \subseteq
    \bigcup_{i=1}^{\infty} G_i \stand G_i \in \mathcal{E} \subseteq
    \set{\text{open sets}}}
  \]
  then $\mu$ is inner regular and outer regular.
\end{proposition}
\begin{remark}
  This applies to Lebesgue measure and Lebesgue--Stieltjes measure for $n = 1$.
\end{remark}

\begin{corollary}
  For a measure space $(\R^n, \A_{\mu^*},\mu)$ as described in
  \Cref{prp:implies-regular} (think about $(\R^n,\L(\R^n),m)$)
  the following are equivalent:
  \begin{enumerate}
    \item[(0)] $E \in \A_{\mu^*}$
    \item[(1)]
      \begin{enumerate}
        \item[(a)] For all $\epsilon > 0$ there exists an open set
          $G \supseteq E$ such that $\mu^*(G \setminus E) \le \epsilon$.
        \item[(b)] There exists $V \supseteq E$ such that $V \in G_{\delta}$
          and $\mu(V \setminus E) = 0$.
      \end{enumerate}
    \item[(2)]
      \begin{enumerate}
        \item[(a)] For all $\epsilon > 0$ there exists a closed
          \footnote{If $\mu^*(E) < \infty$ we can choose a compact $F$.}
          set
          $F \subseteq E$ such that $\mu^*(E \setminus F) \le \epsilon$.
        \item[(b)] There exists $H \subseteq E$ such that $H \in F_{\sigma}$
          and $\mu(E \setminus H) = 0$.
      \end{enumerate}
    \item[(3)]
      \begin{enumerate}
        \item[(a)] For all $\epsilon > 0$ there exist an open $G$ and a closed
          $F$ such that $F \subseteq E \subseteq G$ such that
          $\mu(G \setminus F) \le \epsilon$.
        \item[(b)] There exist
          $F_{\delta} \ni H \subseteq E \subseteq V \in G_{\delta}$ such that
          $\mu(V \setminus H) = 0$.
      \end{enumerate}
    \item[(4)] For $n = 1$ if $\mu^*(E) < \infty$, then for all $\epsilon > 0$
      there exists $U = \bigcup_{i=1}^{N} (a_i,b_i)$ such that
      $\mu^*(E \triangle U) \le \epsilon$.
  \end{enumerate}
\end{corollary}
\begin{proof}
  To be added
\end{proof}

\begin{corollary}
  $\L(\R^n)$ is the minimal completion of $B(\R^n)$ with regards to
  Lebesgue negligible sets.
\end{corollary}
\begin{proof}
  To be added
\end{proof}

\begin{proposition}
  Lebesgue measure is invariant under translations.
  This means that $m^*(A) = m^*(A + h)$ for all $A \subseteq \R^n$
  and $h \subseteq \R^n$.
  It follows that $A$ is Lebesgure measurable if and only if $A + h$ is
  Lebesgure measurable.
\end{proposition}
\begin{proof}
  To be added
\end{proof}

\begin{proposition}
  Similarly $m^*(x \cdot A) = \abs{c}^n \cdot m^*(A)$ for all $c \in \R$.
  It follows that $A$ is Lebesgure measurable if and only if $c\cdot A$ is
  Lebesgure measurable for $c \neq 0$.
\end{proposition}

\begin{proposition}
  For every $T \in \GL(n)$ we have that $m^*(TA) = \abs{\det(T)} m^*(A)$.
\end{proposition}

\begin{proposition}
  There exists a set $P \subseteq [0,1]$ that is not Lebesgure measurable.
\end{proposition}
\begin{proof}
  Define an equivalence relation on $\R$
  \[
    x \sim y \iff x - y \in \Q.
  \]
  We see that
  \[
    \abs{\R / \sim} = \abs{\set{[x_0] = x_0 + \Q \mid x_0 \in \R}} > \aleph_0.
  \]
  Using the \textit{axiom of choice} we can take choose a representative in
  $[0,1]$ from each equivalency class,
  and set $P \subseteq [0,1]$ to be the set of all representatives.
  By definition we see that
  \[
    P + \Q =
    \bigsqcup_{p \in P} p + \Q =
    \bigsqcup_{p \in P} [p] = \R
  \]
  and since $[0,1]^c = [0,1] + [-1,1]^c \supseteq P + \Q \cap [-1,1]^c$,
  we see that
  \[
    [0,1] \subseteq
    P + \del{\Q \cap [-1,1]} \subseteq
    [0,1] + [-1,1] \subseteq
    [-1,2].
  \]
  Thus if we set $\set{q_i}_{i=1}^{\infty} = \Q \cap [-1,1]$ and define
  $P_i = q_i + P$ we will get
  \[
    [0,1] \subseteq \bigsqcup_{i=1}^{\infty} P_i \subseteq [-1,2].
  \]
  Since Lebesgue measure is invariant under translations we get that
  for all $i \geq 1$ that $m(P_i) = m(P)$ and then
  \[
    m\del{[0,1]} = 1 \le
    \sum_{i=1}^{\infty} m(P_i) =
    \sum_{i=1}^{\infty} m(P) \le
    3 = m\del{[-1,2]}.
  \]
  However this is a contradiction since there does not exist a number
  $m(P) \in [0,\infty]$ such that $1 \le \sum_{i=1}^{\infty} m(P) \le 3$
  which completes the proof.
\end{proof}

\begin{remark}
  Recall that we defined the Cantor set as
  \[
    C := \set{x \in [0,1] \middle\vert x =
      \sum_{n=1}^{\infty} \frac{a_n}{3^n} \st a_n \in \set{0,2}}.
  \]
\end{remark}

\begin{proposition}
  The Cantor set satisfies the following properties:
  \begin{enumerate}
    \item[(1)] It is compact;
    \item[(2)] We have $m(C) = 0$;
    \item[(3)] $|C| = \aleph = 2^{\aleph_0}$.
  \end{enumerate}
\end{proposition}
\begin{proof} \phantom{}
\begin{enumerate}
\item[(1)] We notice that
  $C = \bigcap_{n} K_n$ such that
  \[
    K_n := \bigsqcup
    \set{x_0 + [0,3^{-n}] \middle\vert x_0 = 
    \sum_{n=1}^{\infty} \frac{a_n}{3^n} \st a_n \in \set{0,2}}.
  \]
  Since $K_n$ is a union of a finite amount of compact sets it is compact,
  and since $C$ is an intersection of a countable amount of compact sets
  it is compact.
\item[(2)] We have from the monotonicity of $m$ that:
  \[
    0 \le m(C) \le m(K_n) = \frac{2^n}{3^n} \taking{n \to \infty} 0.
  \]
\item[(3)] Let $x \in C$. Then it has a unique ternary expansion without
  the digit $1$:
  \[
    x = \sum_{n=1}^{\infty} \frac{a_n}{3^n} \st a_n \in \set{0,2}
  \]
  so it is clear that
  \[
    |C| = \abs{\set{0,2}^{\N}} = 2^{\aleph_0} = \aleph = \abs{[0,1]}.
  \]
\end{enumerate}
\end{proof}

\begin{definition}[Devil's staircase]
  We can construct the almost-bijection $f \colon [0,1] \to C$ explicitly.
  Define $b_n = \frac{a_n}{2} \in \set{0,1}$ and define
  \[
      y = f(x) := \sum_{n=1}^{\infty} \frac{b_n}{2^n}.
  \]
  It is clear that $f(C) = [0,1]$ because every number in $[0,1]$ has a binary
  expansion, and $f$ is injective except for the numbers that have two
  binary expansions.
  We notice that $f$ is monotonic (nondecreasing) because $x < y$
  implies $f(x) \le f(y)$ and if $x$, $y$ are points
  on the boundary of a $G_n$ set we even have $f(x) < f(y)$.
  We can expand the function such that $f \colon [0,1] \to [0,1]$ by
  defining it to be constant on $G_n$, and this is well defined because
  $f(x) = f(y)$ for $x,y \in \partial G_n$.
  Since $f$ is monotonic it can only have a jump discontinuities,
  but since $f$ is onto $[0,1]$ it doesn't have such points and thus continuous.
  The function $f$ is called Cantor's function
  (or sometimes the Devil's staircase).
  It is continuous everywhere and has zero derivative almost everywhere,
  but value still goes from $0$ to $1$ as its argument reaches from $0$ to $1$.
  It is a very important counterexample in analysis.
\end{definition}

\begin{proposition}
  There exists a Lebesgue measurable function that is not a Borel set.
  That is,
  \[
    \L(\R) \setminus B(\R) \neq \emptyset.
  \]
\end{proposition}
\begin{proof}
  We notice that
  \[
    |B(\R)| = \aleph_1 = 2^{\aleph},
  \]
  but we have that
  \[
    |\L(\R)| \geq 2^{\aleph_1} > \aleph_1
  \]
  because every subset of $C$ is of measure $0$, (and since the measure
  is complete) it is measurable, and there are $2^{|C|}$ such subsets.
\end{proof}
\begin{proof}
  Let $f$ be the Cantor function, and define
  \[
    g(y) = \inf\set{x \in [0,1] \mid f(x) = y}
  \]
  as its inverse function.
  The function $g$ is strictly increasing, and thus injective, 
  but it is not continuous.
\end{proof}

\newpage

\section{Measurable functions}
\subsection{Definitions}
\begin{definition}[Inverse image]
  Given a map $f \colon X \to Y$.
  The inverse image map $f^{-1} \colon 2^{Y} \to 2^{X}$ is given by
  \[
    f^{-1}(E) = \set{x \in X \mid f(x) \in E}, \quad E \subseteq Y.
  \]
\end{definition}
\begin{remark}
  The inverse image preserves taking complements, unions, and intersections
  of any cardinality.
\end{remark}

From this remark we get the following proposition.

\begin{proposition}
  \label{prp:inverse-functions-induces-sigma-algebra}
  Let $\A_Y$ be a $\sigma$-algebra on $Y$.
  Then $f^{-1}(\A_Y) = \set{f^{-1}(E) \mid E \in \A_Y} \subseteq 2^X$
  is a $\sigma$-algebra on $X$.
\end{proposition}

\begin{definition}[Measurable function]
  Let $(X, \A_X)$ and $(Y, \A_Y)$ be measurable spaces,
  and let $f \colon X \to Y$ be a function.
  Then $f$ is called a measurable function if for all $E \in \A_Y$ we
  have $f^{-1}(E) \in \A_X$.
  That is $f^{-1}(\A_Y) \subseteq \A_X$.
\end{definition}

\begin{remark}
  A composition of measurable functions is measurable.
  This follows from the rule $(g \circ f)^{-1} = f^{-1} \circ g^{-1}$.
\end{remark}

\begin{lemma}
  Let $\A_Y = \sigma(\mathcal E)$ such that $\mathcal E \subseteq 2^{Y}$.
  Then $f \colon X \to Y$ is measurable if and only if for all
  $E \in \mathcal E$ we have $f^{-1}(E) \in \A_X$.
\end{lemma}
\begin{proof}
  The first direction follows immediately from the definitions.
  Assume that for all $E \in \mathcal E$ we have $f^{-1}(E) \in \A_X$.
  We notice that the set
  \[
    \set{E \subseteq Y \mid f^{-1}(E) \in \A_X}
  \]
  is a $\sigma$-algebra from \Cref{prp:inverse-functions-induces-sigma-algebra},
  and it is given by assumption that it contains $\mathcal E$.
  This implies that it contains $\A_Y = \sigma(\mathcal E)$ which completes
  the proof.
\end{proof}

\begin{corollary}
  Let $X$, $Y$ be topological spaces.
  Then a continuous funtion $f \colon (X,B(X)) \to (Y,B(Y))$ is measurable.
\end{corollary}
\begin{proof}
  We have that $B(Y) = \sigma\del{\set{\text{open sets}}}$,
  but since for every open $G \subseteq Y$ we have that $f^{-1}(G) \subseteq X$
  is open, we get that $f^{-1}(G) \in B(X)$ which completes the proof.
\end{proof}

\begin{remark}
  Notice that measurable functions are a lot more general than continuous
  functions because we require $f^{-1}(G)$ to be Borel not necessarily open.
\end{remark}

\begin{remark}
  If $(X,\A)$ is a measurable space and $f \colon X \to \R$,
  we always endow $\R$ with $B(\R)$.
\end{remark}

\begin{corollary}
  The following are equivalent:
  \begin{enumerate}
    \item[(1)] $f \colon (X,\A) \to \R$ is measurable.
    \item[(2)] For all $a \in \R$ we have $f^{-1}\del{[a,\infty)} \in \A$.
    \item[(3)] For all $a \in \R$ we have $f^{-1}\del{(a,\infty)} \in \A$.
    \item[(4)] For all $b \in \R$ we have $f^{-1}\del{(\ninfty,b]} \in \A$.
    \item[(5)] For all $b \in \R$ we have $f^{-1}\del{(\ninfty,b]} \in \A$.
  \end{enumerate}
  Where for example $f^{-1}\del{(\ninfty,b)} = \set{x \in X \mid f(x) < b}$.
\end{corollary}

\begin{definition}[Borel/Lebesgue measurable function]
  A function $f \colon \R^n \to \R$ is said to be Borel or Lebesgue measurable
  (with respect to its domain $\R^n$) if
  \[
    f \colon
    \begin{array}{c}
      (\R^n,B(\R^n)) \\
      (\R^n,\L(\R^n))
    \end{array}
    \to (\R, B(\R))
  \]
  is measurable.
\end{definition}
\begin{remark}
  We always consider the $\sigma$-algebra on the domain of the function,
  not the range.
\end{remark}

\begin{remark}
  A continuous function $f \colon \R^n \to \R$ is Borel (and thus Lebesgue)
  measurable, and a function $g$ that is equal to $f$ almost everywhere,
  is also Lebesgue measurable (but not necessarily Borel measurable).
  That is because if
  \[
    m\del{\underbrace{\set{x \in \R^n \mid f(x) \neq g(x)}}_{N_0}} = 0
  \]
  then for any $E \in B(\R)$ we have that $g^{-1}(E) = f^{-1}(E) \triangle N$
  for some negligible subset $N$ of $N_0$.
\end{remark}

\begin{example}[Lebesgure measurable but not Borel measurable]
  The function $f = 0$ is Borel and Lebesgure measurable,
  but for $E$ that is Lebesgure measurable and not Borel measurable,
  such that $m(E) = 0$
  the function $g = \mathbbm{1}_{E}$ is not Borel measurable even though
  it is equal to $f$ almost everywhere.
\end{example}

\begin{remark}
  If $f,g \colon \R \to \R$ such that $f$ is Borel measurable,
  $g$ is Borel/Lebesgue measurable,
  then $f \circ g$ is Borel/Lebesgue measurable
  (because $(g \circ f)^{-1} = f^{-1} \circ g^{-1}$).
  However, if $f$ is Lebesgue measurable,
  then $f \circ g$ might not be measurable at all.
\end{remark}

\begin{definition}[Extended real number line]
  The extended real number line $\overline \R$ is defined as $[\ninfty,+\infty]$
  or $\R \cup \set{\ninfty,+\infty}$.
  It can be turned into a totally ordered set by defining
  $\ninfty \le a \le +\infty$ for all $a \in \overline \R$.
  Since it is totally ordered, we can endow it with the order topology,
  which is the topology generated by the subbase of open rays.
  In this way, the space $\overline \R$ is compact,
  every subset of $\overline \R$ has a supremum and an infiumum
  (the infimum of the empty set is $+\infty$, and its supremum is $\ninfty$.
  With this topology $\overline \R$ is also homeomorphic to the unit interval
  $[0,1]$.
  Thus, the topology is metrizible, correspoinding to the standard metric
  on the interval.
  However, there is not metric that is an extension of the standard metric on
  $\R$.
\end{definition}
\begin{remark}
  This implies that $B(\overline \R)$ is defined as the all the sets
  of the form $B \in B(\R)$ or $B \cup \set{\infty}$ or $B \cup \set{\ninfty}$
  or $B \cup \set{\ninfty,\infty}$.
\end{remark}
\begin{remark}
  A function $f \colon (X,\A) \to \overline \R$ is said to be measurable
  if $f \colon (X,\A) \to (\overline \R, B(\overline \R))$ is measurable.

  A function $f \colon (X,\A) \to \C$ is said to be measurable
  if $f \colon (X,\A) \to (\C, B(\C))$ is measurable.

  Since we have that $\C \simeq \R^2$ we have that
  \[
    B(\C) = B(\R^2) = B(\R) \otimes B(\R).
  \]
  This implies that $f$ is measurable if
  \[
    \Im(f) \colon (X,\A) \to (\R,B(\R)) \tand
    \Re(f) \colon (X,\A) \to (\R,B(\R))
  \]
  are both measurable because
  \[
    (\Re f)^{-1}(A) \cap (\Im f)^{-1}(B) = f^{-1}(A \times B).
  \]
\end{remark}

\begin{proposition}
  \label{prp:measurable-functions-arithmetic}
  Let $f,g \colon (X,\A) \to \R$ be measurable, and let $c \in \R$.
  Then functions of the form
  \begin{enumerate}
    \item[(1)] $c \cdot f$,
    \item[(2)] $f + g$,
    \item[(3)] $f \cdot g$,
    \item[(4)] $\frac{f}{g}$ (if $g \neq 0$),
  \end{enumerate}
  are all measurable as well.
\end{proposition}
\begin{proof}
  To be added.
\end{proof}
\begin{remark}
  We have that \Cref{prp:measurable-functions-arithmetic} is true
  for $f,g \colon (X,\A) \to \overline \R$ as well when declaring
  \[
    (\pinfty) + (\pinfty) = \pinfty \quad
    (\ninfty) - (\ninfty) = \ninfty \quad
    0 \cdot (\pm\infty) = 0 \quad
    (\pinfty) + (\ninfty) = a \in \overline \R
  \]
  for some fixed $a \in \overline \R$.
\end{remark}

\begin{proposition}
  Let $f,g \colon (X,\A) \to \overline \R$ be measurable.
  The the functions $\abs{f}$, $\max(f,g)$, and $\min(f,g)$ are measurable.
\end{proposition}
\begin{proof} \phantom{}
\begin{enumerate}
  \item[(1)] This is simply a composition with $x \mapsto \abs{x}$.
  \item[(2 + 3)] We can see this by noticing that
  \begin{align*}
    \set{\max(f,g) < b} &= \set{f < b} \cap \set{g < b} \\
    \set{\min(f,g) < b} &= \set{f < b} \cup \set{g < b}.
  \end{align*}
  or by ??
  % add later.
\end{enumerate}
\end{proof}

\subsection{Pointwise convergence}
\begin{proposition}
  \label{prp:more-measurable-functions}
  Let $\set{f_i \colon (X,\A) \to \overline \R}_{i=1}^{\infty}$ be a sequence
  of measurable functions.
  Then the functions
  \[
    \limsup_{i \to \infty} f_i \quad
    \liminf_{i \to \infty} f_i \quad
    \sup_{i \to \infty} f_i \quad
    \inf_{i \to \infty} f_i
  \]
  are all measurable.
\end{proposition}
\begin{remark}
  All the limits are pointswise.
  That is, $(\sup_{i} f_i)(x) = \sup_{i}(f_i(x))$.
\end{remark}
\begin{proof}
  To be added (not hard).
\end{proof}

\begin{corollary}
  Let $\set{f_i \colon (X,\A) \to \overline \R}_{i=1}^{\infty}$ be a sequence
  of measurable functions that converge pointwise to $f$.
  Then $f \colon (X,\A) \to \overline \R$ is measurable.
\end{corollary}
\begin{proof}
  \[
    \lim_{i \to \infty} f_i =
    \underbrace{\liminf_{i \to \infty} f_i}_{\text{measurable}}.
  \]
\end{proof}

\subsection{Simple functions}
\begin{definition}[Indicator function]
  An indicator function of a set $E \subseteq X$ is defined as
  \[
    \mathbbm{1}_E(x) =
    \chi_E(x) =
    \begin{cases}
      1, &x \in E \\
      0, &\text{otherwise}
    \end{cases}.
  \]
\end{definition}
\begin{remark}
  $\mathbbm 1_E$ is measurable if and only if $E$ is measurable.
  That is because
  \[
    \set{\mathbbm{1}_{E} \geq b} =
    \begin{cases}
      \emptyset, &b > 1 \\
      E, &0 < b \le 1 \\
      X, &b \le 0
    \end{cases}.
  \]
\end{remark}

\begin{definition}[Simple function]
  \label{def:simple-function}
  A function $\phi \colon (X,\A) \to \R$ (not $\overline\R$)
  is called simple if
  \[
    \phi(X) =
    \sum_{j=1}^{N} c_j \cdot \mathbbm{1}_{E_j}, \quad
    c_j \in \R \stand E_j \in \A.
  \]
\end{definition}
\begin{remark}
  We see that $\phi$ is a simple function if and only if it is measurable
  and $\range(\phi)$ is a finite set.
\end{remark}
\begin{remark}
  There could be more than one representation of the same simple function
  as seen in \Cref{def:simple-function}.
  The representation is called standard if $\set{c_i}_{i=1}^{\infty}$ are all
  differenet and $E_j$ are disjoint.
  This representation is unique because we can set
  $\set{c_i} = \range(\phi)$ and $E_j = \phi^{-1}\del{\set{c_i}}$.
\end{remark}
\begin{remark}
  The sum and product of simple functions is simple.
\end{remark}

\begin{proposition}
  Let $f \colon (X,\A) \to [0,\infty]$ be a nonnegative measurable function.
  % nonnegative is here for emphasis
  Then there exists a sequence of simple functions $\set{\phi_n}$ such that
  \[
    0 \le \phi_1 \le \phi_2 \le \cdots
  \]
  is a pointwise increasing function and $\phi_n \nearrow f$ pointwise.
  Moreover, $\phi_n \taking{n \to \infty} f$ uniformly on any set $E$
  such that $f$ is bounded on $E$.
\end{proposition}
\begin{proof}
  Given $n \in \N$, we split the interval $(0,2^n]$ to $2^{2n}$ subintervals
  of length $2^{-n}$ as follows
  \[
    I_{k,n} = (k \cdot 2^{-n}, (k+1) \cdot 2^{-n}] \quad
    k = 0,1,2,\dots,2^{2n}-1.
  \]
  and set $J_n = (2^n,\infty)$.
  We also define
  \[
    F_n = f^{-1}(J_n) \tand
    E_{n,k} = f^{-1}(I_{n,k}) \tand
    \phi_n =
    \sum_{k=0}^{2^{2n}-1}
    k \cdot 2^{-n} \cdot \mathbbm{1}_{E_{k,n}} + 2^n \cdot \mathbbm{1}_{F_n}.
  \]
  % Add sketch?
  The sequence $\phi_n$ is increasing because we just refine the partitions.
  We also have that if $f\vert_{E} \le 2^{n_0}$ then for all $n_0 \le n$
  we have that $0 \le f \cdot \phi_n \le 2^{-n}$ on $E$.
  This shows that $\phi_n \taking{n \to \infty} f$ uniformly on any set $E$
  such that $f$ is bounded on $E$.
  In particular we have a pointwise convergence for any $x$ such that
  $f(x) < \infty$.
  If $f(x) = \infty$ then the approximation of $\phi_n$ is given by
  $2^n \taking{n \to \infty} \infty$ so we have that $\phi_n \nearrow f$
  pointwise which completes the proof.
\end{proof}

\subsection{Almost everywhere}
Recall the following definion.
\begin{definition}[Almost everywhere]
  Let $(X,\A,\mu)$ be a measure space.
  We say that a certain property is true \textit{almost everywhere} if it is
  true for any $x \in X \setminus E$ such that $E$ is of measure zero
  ($\mu$-negligible).
\end{definition}
\begin{remark}
  It is possible that $\overline B = \set{x \mid \neg P(x)} \notin \A$.
  We only require that $\overline B \subseteq B \in \A$ such that $\mu(B) = 0$.
  Altenratively, that $\overline \mu(\overline B) = 0$ where
  $(X,\overline \A,\overline \mu)$ is the completion of $(X,\A,\mu)$.
\end{remark}

\begin{proposition}
  Let $(X,\A,\mu)$ be a complete measure space,
  and let $f,g \colon X \to \overline \R$.
  Then if $f$ is measurable,
  and $f = g$ almost everywhere ($\mu(\underbrace{f \neq g}_{N}) = 0$),
  then $g$ is measurable.
\end{proposition}
\begin{proof}
  From the completeness of the space $(X,\A,\mu)$ we get that
  \[
    \set{g < b} =
    \left[
    \set{f < b} \cap N^c
    \right] \sqcup
    \left[
    \set{g < b} \cap N
    \right] \in \A.
  \]
\end{proof}

\begin{proposition}
  Let $(X,\A,\mu)$ is a complete measure space,
  let $\set{f_n \colon (X,\A) \to \overline \R}_{i=1}^{\infty}$ be a sequence
  of measurable function,
  and let $f_n \taking{n \to \infty} f$ almost everywhere,
  that is
  \[
    \mu\del{\set{x \in X \mid
        \nexists \lim_{n \to \infty} f_n(x)
        \textnormal{ or }
        \exists \lim_{n \to \infty} f_n(x) \neq x}} = 0.
  \]
  Then $f$ is measurable.
\end{proposition}
\begin{proof}
  Define $g = \limsup_{n \to \infty} f_n$.
  It is measurable from \Cref{prp:more-measurable-functions}.
  Since $f = g$ $\mu$-almost everywhere, from completeness $f$ is measurable.
\end{proof}

\begin{proposition}
  \label{prp:overline-a-e-noverline}
  Let $(X,\overline \A,\overline \mu)$ be the completion of the measure
  space $(X,\A,\mu)$.
  Let $\overline f \colon X \to \overline \R$ be $\overline \A$-measurable.
  Then there exists a $\A$-measurable function $f \colon X \to \R$ such
  that $\overline f = f$ almost everywhere.
\end{proposition}

\begin{proof}
  To be added.
\end{proof}

\begin{corollary}
  Let $f_n \taking{n \to \infty} \overline f$ $\mu$-almost everywhere.
  Then it is possible that $\overline f$ is not $\A$-measurable
  (if $(X,\A,\mu)$ is not complete) but there exists a measurable
  function $f$ such that $\overline f = f$ almost everywhere,
  and thus $f_n \taking{n \to \infty} f$ $\mu$-almost everywhere.
\end{corollary}

\newpage

\section{Integration}
\subsection{Integrals of simple functions}

Set a measure space $(X,\A,\mu)$ and denote
$\L^+ := \set{f \colon X \to [0,\infty] \mid \text{$f$ is measurable}}$.
\begin{remark}
  If $\phi \in \L^+$ is simple, then
  \[
    \phi = 
    \sum_{j=1}^{N} c_j \cdot \mathbbm{1}_{E_j}, \quad
    c_j \in [0,\infty] \stand E_j \in \A.
  \]
\end{remark}

\begin{definition}[Integration of simple functions]
  Let $\phi \in \L^+$ be simple.
  Then we define
  \[
    \int \phi \dmu = 
    \sum_{j=1}^{N} c_j \cdot \mu(E_j) \in [0,\infty].
  \]
\end{definition}

\begin{remark}
  We can verify that this definition is not dependent on the representation
  of $\phi$.
  Thus we usually compute this integral using the simple function's
  standard representation.
\end{remark}
\begin{remark}
  The convention is that
  $\underset{=c_j}{0} \cdot \underset{=\mu(E_j)}{\infty} = 0$.
\end{remark}
\begin{remark}
  We denote
  \[
    \int \phi =
    \int \phi \dmu =
    \int \phi(x) \dmu(x) =
    \int \phi(x) \mu(\dx)
  \]
  and for all $A \in \A$ 
  \[
    \int_{A} \phi \dmu =
    \int \underbrace{\phi \cdot \mathbbm{1}_{A}}_{\text{simple}} \dmu.
  \]
\end{remark}

\begin{proposition}
  Let $\phi, \psi \in \L^+$ be simple functions. Then
  \begin{enumerate}
    \item[(1)] For all $a \geq 0$ we have $\int a \phi = a \int \phi$.
    \item[(2)] $\int (\phi + \psi) = \int \phi + \int \psi$.
    \item[(3)] If $\phi \le \psi$ then $\int \phi \le \int \psi$.
  \end{enumerate}
\end{proposition}
\begin{proof} \phantom{}
\begin{enumerate}
  \item[(1)] Immediate result.
  \item[(2)] 
  \item[(2)] 
\end{enumerate}
\end{proof}

\begin{definition}[Integrability]
  Let $f \in \L^+$.
  We define
  \[
    \int f \dmu :=
    \sup\set{\int \phi \dmu \middle\vert 0 \le \phi \le f \stand \phi 
    \text{ is simple}}.
  \]
  A function $f \in \L^+$ is called integrable with respect to $\mu$ if
  $\int f \dmu < \infty$.
  We define
  \[
    \int_{A} f \dmu = \int f \cdot \mathbbm{1}_{A} \dmu.
  \]
\end{definition}
\begin{remark}
  This definition generalized the previous definition for $f = \phi$.
  This follows from property $(3)$ in the previous proposition.
\end{remark}
\begin{remark}
  One may ask why don't we approximate the integral from above?
  The answer is that while for bounded functions with $|A| < \infty$
  we can equivalently define
  \[
    \int f \dmu :=
    \inf\set{\int \phi \dmu \middle\vert 0 \le f \le \phi \stand \phi
    \text{ is simple}}.
  \]
  However, the situation is not so simple if the function is unbounded
  or $|A| = \infty$. Here are a few examples:
  \begin{itemize}
    \item The function $f(x) = x^{-2}$ is integrable over $A = (1,\infty)$.
      However, for any simple function such that $f \le \phi$ we have
      $\int_{A} \phi \dmu = \infty$ which means it is not integrable.
    \item The function $f(x) = x^{-0.5}$ is integrable over $A = (0,1)$.
      However, there's no simple function $s$ such that $f \le s$ and
      $s$ is finite almost everywhere.
  \end{itemize}
\end{remark}

\begin{proposition}
  \label{prp:l-plus-properties}
  Let $\psi, \phi, \phi_1, \phi_2 \in \L^+$.
  Then
  \begin{enumerate}
    \item[(1)] For all $a \geq 0$ we have $\int a \phi = a \int \phi$.
    \item[(2)] $\int (\phi_1 + \phi_2) = \int \phi_1 + \int \phi_2$.
    \item[(3)] If $0 \le \psi \le \phi$ then $0 \le \int \psi \le \int \phi$.
  \end{enumerate}
\end{proposition}

\subsection{Monotone convergence theorem}
\begin{theorem}[Monotone convergence theorem]
  Let $\set{f_n} \subseteq \L^+$ be a nondecreasing sequence of functions
  that converges pointwise to $f$.
  Then
  \[
    \lim_{n \to \infty} \int f_n \dmu = \int f \dmu.
  \]
\end{theorem}
\begin{remark}
  The function $f$ is measurable because $f = \sup_{n} f_n$.
\end{remark}
\begin{remark}
  The theorem is false for the Riemann integral.
  Since $\Q$ is countable we can construct a well order on it such that
  $\set{q_k}_{k=1}^{\infty} = \Q$ and then define
  \[
    f_n(x) = \begin{cases}
      1, &x\in \bigcup_{k=1}^{n} q_k \\
      0, &\text{otherwise}
    \end{cases}
  \]
  but the limit function is $f = \mathbbm{1}_{\Q \cap [0,1]}$ which is not
  Riemann integrable.
\end{remark}
\begin{proof}
  From property $(3)$ of \Cref{prp:l-plus-properties}
  we have
  \[
    \int f_1 \le \int f_2 \le \cdots \le \int f_n \le \int f
  \]
  so the sequence $\int f_n$ converges and we have
  \[
    \lim_{n \to \infty} \int f_n \dmu \le \int f \dmu.
  \]
  Now set $\alpha \in (0,1)$ and let $\phi$ be a simple function such that
  $0 \le \phi \le f$.
  We define
  \[
    A_n =\set{x \in X \colon f_n(x) \geq \alpha \cdot \phi(x)}
  \]
  and see that they are measurable, increasing, and $\bigcup A_n = X$.
  Thus
  \[
    \int f_n \geq
    \int f_n \mathbbm{1}_{A} \geq
    \int \alpha \phi \cdot \mathbbm{1}_{A_n} =
    \alpha \int_{A_n} \phi.
  \]
  and then
  \[
    \int_{A_n} \phi =
    \int \phi \mathbbm{1}_{A_n} =
    \sum_{j=1}^{N} c_j \mu\del{E_j \cap A_n} %%%% weird up arrow
    \sum_{j=1}^{N} c_j \mu(E_j) = \int \phi.
  \]
  Thus finally we have
  \[
    \lim_{n \to \infty} \int f_n \geq \alpha \int \phi
  \]
  for an arbitrary $\alpha \in (0,1)$ so we can take $\alpha \to 1$
  and get
  \[
    \lim_{n \to \infty} \int f_n \geq \int \phi.
  \]
  We can now take $\sup_{\phi}$ and get that
  \[
    \lim_{n \to \infty} \int f_n \geq \int f
  \]
  which completes the proof.
\end{proof}
\begin{corollary}
  An equivalent definition to the integral of $f \in \L^+$ is
  \[
    \int f \dmu = \lim_{n \to \infty} \int \phi_n \dmu
  \]
  for any sequence of positive simple functions $0 \le \phi_n \to f$.
  This is because we have shown that there exists such a sequence and the
  limit is not dependent on the choice of sequence.
\end{corollary}

\begin{corollary}
  For $f,g \in \L^+$ we have that
  \[
    \int f + g =
    \int f + \int g.
  \]
\end{corollary}
\begin{proof}
  There exists sequences of simple functions $\phi_n \nearrow f$ and 
  $\psi_n \nearrow g$, and thus $\phi_n + \psi_n \nearrow f + g$
  and then from MCT we have
  \[
    \int f + g =
    \lim_{n \to \infty}\del{\int \phi_n + \int \psi_n} =
    \int f + \int g.
  \]
\end{proof}

\begin{corollary}[Switching integrals and sums]
  For $\set{f_n}_{n=1}^{\infty} \subseteq \L^+$ we have
  \[
    \int \sum_{n=1}^{\infty} f_n =
    \sum_{n=1}^{\infty} \int f_n.
  \]
\end{corollary}
\begin{proof}
  We can show this is true for a finite $n$ by induction on $g$,
  and in the general case we have
  \[
    \int \sum_{n=1}^{\infty} f_n =
    \int \lim_{N \to \infty} \sum_{n=1}^{N} f_n 
    \underset{\text{MCT}}{=}
    \lim_{N \to \infty} \int \sum_{n=1}^{N} f_n =
    \lim_{N \to \infty} \sum_{n=1}^{N} \int f_n =
    \sum_{n=1}^{\infty} \int f_n
  \]
\end{proof}

\begin{remark}
  Recall that in general, if $f_n \taking{\text{pointwise}} f$ we don't
  necessarily have that
  \[
    \lim_{n \to \infty} \int f_n \dmu = \int f \dmu.
  \]
  The limit might not even exist.
  For example, we can choose $f_n = \mathbbm{1}_{[n,n+1)}$ with $\mu = m$
  being Lebesgue measure.
  We have that $f_n \taking{n \to \infty} 0$ pointwise,
  but for all $n \geq 1$ we have $\inf f_n \dm = 1$ and thus
  \[
    1 = \lim_{n \to \infty} \int f_n \dm \neq
    \int 0 \dm = 0.
  \]
\end{remark}

% Another example + This is because the 'mass' is going to \infty?

\begin{remark}
  The following lemma shows us that even if the monotonic convergence theorem
  does not hold for non monotonic function sequences,
  at least one part of it always holds.
\end{remark}

\begin{lemma}[Fatou's lemma]
  Let $\set{f_n}_{n=1}^{\infty} \subseteq \L^+$.
  Then
  \[
    \int \liminf_{n \to \infty} f_n \le
    \liminf_{n \to \infty} \int f_n.
  \]
\end{lemma}

\begin{remark}
  The direction of the inequality is affected by
  $\set{f_n}_{n=1}^{\infty} \subseteq \L^+$.
  For nonpositive functions we have 
  \[
    \int \limsup_{n \to \infty} f_n \geq
    \limsup_{n \to \infty} \int f_n.
  \]
\end{remark}
\begin{proof}
  Set $k$, and then for all $j \geq k$ we have that that
  \[
    \inf_{n \geq k} f_n \le f_j \implies
    \int \inf_{n \geq k} f_n \le \int f_j \implies
    \int \inf_{n \geq k} f_n \le \inf_{j \geq k} \int f_j \implies
  \]
  TO BE CONTINUED
\end{proof}

% Go over this proof
\begin{proposition}
  Let $f \in \L^+$.
  Then $\int f \dmu = 0$ if and only if $f = 0$ $\mu$-almost everywhere.
\end{proposition}
\begin{proof}
  First suppose that $\phi$ is simple.
  \[
    \phi= \sum_{j=1}^{N} c_j \cdot \mathbbm{1}_{E_j}.
  \]
  This case is clear because
  \begin{align*}
    0 =
    \int \phi \dmu =
    \sum_{j=1}^{N} c_j \cdot \mu(E_j) &\iff
    \forall j \quad c_j \mu(E_j) = 0 \\ &\iff
    c_j = 0 \stor \mu(E_j) = 0 \iff
    \phi = 0 \text{ $\mu$-almost everywhere}.
  \end{align*}
  Now for a general function, first assume that $f = 0$ almost everywhere.
  Then if $0 \le \phi \le f$ is simple then $\phi = 0$ almost everywhere
  and thus $\int \phi \dmu = 0$ so $\int f = \sup \int \phi \dmu = 0$.
  
  For the second direction we need to use Markov's inequality.
  \begin{lemma}[Markov's inequality]
    Let $f \in \L^+$. Then for all $r < 0$ we have
    \[
      \mu(f \geq r) \le \frac{\int f \dmu}{r}.
    \]
  \end{lemma}
  \begin{proof}
    \[
      r \cdot \mu(f \geq r) =
      \int r \cdot \mathbbm{1}_{\set{f \geq r}} \dmu \le
      \int f \dmu.
    \]
  \end{proof}
  Define $A_n = \set{f \geq \frac{1}{n}}$.
  Assume by contradiction that we don't have $f = 0$ almost everywhere,
  then there would exist $n$ such that $\mu(A_n) > 0$, because otherwise
  we would have
  \[
    \mu(f > 0) = \mu\del{\bigcup A_n}
  \]
  and from continuity from below this is equal to
  $\lim_{n \to \infty} \mu(A_n) = 0$ so $f = 0$ everywhere which is a
  contradiction.
  Thus,
  \[
    f \geq \frac{1}{n} \cdot \mathbbm{1}_{A_n} \implies
    \int f \geq \frac{1}{n} \cdot \mu(A_n) > 0
  \]
  in contradiction, which completes the proof.
\end{proof}

\begin{corollary}
  In the convergence theorems (MCT/Fatou) we can require convergence in
  $\mu$-almost everywhere instead of pointwise convergence.
\end{corollary}
%% RECOLLECTIONS?
%% PROOF OF COROLLARY?

\begin{proposition}
  Let $f \in \L^+$ such that $\int f \dmu < \infty$.
  Then $f$ is finite $\mu$-almost everywhere ($\mu\del{\set{f = \infty}} = 0$).
\end{proposition}
\begin{proof}
  The set $E = \set{f = + \infty}$ is measurable,
  and $f \geq t \cdot \mathbbm{1}_{E}$ for all $t$ so
  \[
    \int f \geq t \int \mathbbm{1}_{E} = t \cdot \mu(E).
  \]
  So if $\mu(E) > 0$ we can take $t \to \infty$ and get a contradiction.
\end{proof}
\begin{exercise}
  If $f \in \L^+$ and $\int f \dmu < \infty$ then $\set{f > 0}$ is
  $\sigma$-finite.
\end{exercise}

\subsection{Integrals of general measurable functions}
For each $f \colon X \to [-\infty, \infty]$ we define
\[
  f^+, f^- \colon X \to [0,\infty]
\]
by
\[
  f^+ = \max\set{f, 0} \tand
  f^- = \max\set{-f, 0}.
\]
It is clear that $f = f^+ - f^-$ and $|f| = f^+ + f^-$.
We also have that $f$ is measurable if and only if $f^-$ and $f^+$
are measurable. They are also nonnegative so if $f$ is measurable thism means
they are in $\L^+$ so we can also take their integrals.
\begin{definition}[Integrability]
  Let $f \colon X \to [-\infty, \infty]$ be measurable, and define
  \[
    \int f \dmu =
    \int f^+ \dmu - \int f^- \dmu
  \]
  if at least one of the intergals is finite.
  If both integrals are finite, we say that $f$ is integrable (or Lebesgure
  integrable if $\mu = m$ is Lebesgure measure).
  In the case both integrals are finite, we have
  \[
    \int |f| \dmu =
    \int f^+ \dmu + \int f^- \dmu < \infty.
  \]
  That is, unlike the Riemann integral,
  $f$ is integrable if and only if $f$ is absolutely integrable!
\end{definition}
\begin{example}[Integral exists but the function is not integrable]
  Define $f(x) = \frac{1}{x} \cdot \mathbbm{1}_{[0,1]}$,
  \[
    \int_{0}^{1} f \dm \underset{\text{MCT}}{=}
    \lim_{\epsilon_n \searrow 0} \int_{\epsilon_n}^{1} \frac{1}{x} \dm =
    \lim_{\epsilon_n \searrow 0} \ln x \big\vert_{\epsilon_n}^{1} = \infty.
  \]
  Thus, the integral exists,
  but the function is not integrable.
\end{example}

\begin{definition}[$\tilde{\L}^1$-space]
  \[
    \tilde{\L}^1(\mu) :=
    \set{f \colon X \to [\ninfty,\infty] \middle\vert \int |f| \dmu < \infty \stand
    f \text{ is measurable}}.
  \]
\end{definition}
\begin{proposition}
  The space $\tilde\L^1$ is a linear space over $\R$, and the integral
  is a linear functional on it.
  Let $f,g \in \tilde\L^1(\mu)$. Then
  \begin{enumerate}
    \item[(1)] $\int c f \dmu = c \int f \dmu$.
    \item[(2)] $\int f + g = \int f + \int g$.
    \item[(3)] $f \le g \implies \int f \le \int g$.
    \item[(4)] $\abs{\int f} \le \int\abs{f}$.
  \end{enumerate}
\end{proposition}
\begin{proof}
  $\tilde\L^1(\mu)$ is a linear space because for $f,g \in \tilde\L^1(\mu)$
  we have
  \[
    \int |a f + b g| \le
    \int |a| |f| + \int |b| |g| \implies
    a f + b g \in \tilde\L^1(\mu)
  \]
  We can now prove the properties as mentioned in the proposition
  \begin{enumerate}
    \item[(1)] It is clear that
      \[
        \int c \cdot f =
        \int (c f)^+ - \int (c f)^- =
        c \del{\int f^+ - \int f^-} =
        c \int f.
      \]
    \item[(2)] Denote $h = f + g$.
    We have that
    \[
      h^+ - h^- =
      f^+ - f^- + g^+ - g^-
    \]
    which implies that
    \begin{align*}
      h^+ + f^- + g^- = f^+ + g^+ + h^- &\implies
      \int h^+ + \int f^- + \int g^- = \int f^+ + \int g^+ + \int h^- 
                                     \\ &\implies
      \int h^+ - \int h^- = \int f^+ - \int f^- + \int g^+ - \int g^-
    \end{align*}
  \item[(3)] Suppose that $g - f \geq 0$, then
    \[
      0 \le \int (f - g) = \int f - \int g
    \]
  \item[(4)] We have that $- |f| \le f \le |f|$ so according to $(3)$ we have
    that
    \[
      - \int |f| \le \int f \le \int |f|
    \]
  \end{enumerate}
\end{proof}
\begin{remark}
  Let $f \colon X \to \C$.
  Then $f$ is integrable if and only if $\Re f$ and $\Im f$ are integrable
  because
  \[
    \int f = \int \Re f + i \int \Im f.
  \]
\end{remark}

\begin{proposition}
  Let $f,g \in \tilde\L^1$. Then
  \begin{enumerate}
    \item[(1)] $\set{f \neq 0}$ is $\sigma$-finite.
    \item[(2)] $f = g$ $\mu$-almost everywhere if and only if
      $\int (f - g) \dmu = 0$ if and only if for all $E \in \A$ we have
      $\int_{E} f \dmu = \int_{E} g \dmu$.
  \end{enumerate}
\end{proposition}
\begin{proof} \phantom{}
\begin{enumerate}
  \item[(1)] We have that $\set{f \neq 0} = \set{|f| \neq 0}$.
    Since $|f| \in \L^+$ we have $\int |f| \dmu < \infty$.
  \item[(2)] $f - g = 0$ $\mu$-almost everywhere if and only if
    $\underset{\in \L^+}{|f - g|} = 0$ $\mu$-almost everywhere if and only
    if $\int |f - g| \dmu = 0$.
\end{enumerate}
\end{proof}

\begin{proposition}
  \label{prp:equal-a-e-implies-equal-integral}
  Let $f,g \in \tilde\L^1(\mu)$.
  Then $f = g$ $\mu$-a.e.\ if and only if
  $\int |f - g| \dmu = 0$ if and only if $\int_{E} f = \int_{E} g$ for all
  $E \in \A$.
\end{proposition}
\begin{proof}
  To be added.
\end{proof}

\begin{definition}[$\L^1(\mu)$ space]
  We define an equivalence relation on $\tilde\L^1(\mu)$ by
  \[
    f \sim g \iff f = g \, (\mu \text{ almost everywhere}).
  \]
  The transitivity follows from the fact that
  $\mu(A \cup B) \le \mu(A) + \mu(B) = 0$.
  We define
  \[
    \L^1(\mu) = \tilde\L^1(\mu) / \sim.
  \]
\end{definition}
\begin{remark}
  The space $\L^1(\mu)$ is a vector space on $\R$.
\end{remark}
\begin{definition}[$\L^1(\mu)$ norm]
  Let $f \in \L^1(\mu)$ be a representative of an equivalence class $[f]$.
  We define
  \[
    \norm{[f]}_{\L^1(\mu)} =
    \norm{f}_{\L^1(\mu)} =
    \int |f| \dmu (< \infty).
  \]
\end{definition}
\begin{remark}
  The norm is well defined by \Cref{prp:equal-a-e-implies-equal-integral}.
\end{remark}
\begin{remark}
  The space $(\L^1(\mu), \norm{\cdot}_{\L^1(\mu)})$ is a complete normed space.
\end{remark}

\begin{remark}
  By \Cref{prp:overline-a-e-noverline} there exists a bijection between
  $\L^1(\mu)$ and $\L^1(\overline \mu)$.
  This means that for all $\overline f$ there exists $f$ such that
  $f \sim \overline f$ and thus we identify them as the same space.
\end{remark}

\subsection{Dominant convergence theorem}

\begin{theorem}[Dominant convergence theorem]
  Let $\set{f_n} \subseteq \L^1(\mu)$ be a sequence such that
  \begin{itemize}
    \item $f_n \to f$ $\mu$-almost everywhere.
    \item There exists an integrable ``dominant'' $g \in \L^1(\mu)$ such that
      $|f_n| \le g$ $\mu$-almost everywhere for all $n$.
  \end{itemize}
  Then $f \in \L^1(\mu)$ and
  \[
    \int f \dmu =
    \lim_{n \to \infty} \int f_n \dmu.
  \]
\end{theorem}
\begin{remark}
  The second requirement is the analog of the finiteness requirement
  in the continuity from above property of measures.
\end{remark}
\begin{proof}
  To be added.
\end{proof}

\begin{proposition}
  Let $\set{f_k} \subseteq \L^1(\mu)$ be a sequence such that
  \[
    \sum_{k=1}^{\infty} \int \abs{f_k} \dmu < \infty.
  \]
  Then
  \begin{enumerate}
    \item[(1)] $\sum_{k=1}^{\infty} \abs{f_k} \in \L^1(\mu)$.
    \item[(2)] There exists $g \in \L^1(\mu)$ such that
      $\sum_{k=1}^{n} f_k \taking{n \to \infty} g$ $\mu$-almost everywhere.
    \item[(3)] $\int g \dmu = \sum_{k=1}^{\infty} \int f_k \dmu$.
  \end{enumerate}
\end{proposition}
\begin{proof} \phantom{}
\begin{enumerate}
  \item[(1)]
  \item[(2)]
  \item[(3)]
\end{enumerate}
\end{proof}

\newpage

\section{Covergence}
\subsection{Types of convergence}
Let $(X,\A,\mu)$ be a measure space.
\begin{description}
  \item[(0) $f_n \taking{n \to \infty} f$ uniformly] if
    \[
      \sup_{x \in X} \abs{f_n(x) - f(x)} \taking{n \to \infty} 0.
    \]
  \item[(1) $f_n \taking{n \to \infty} f$ $\mu$-almost everywhere] if
    \[
      \mu\del{\set{x \in X \mid f_n(x) \not\taking{n \to \infty} f(x)}} = 0.
    \]
  \item[(2) $f_n \taking{n \to \infty} f$ in $\L^1(\mu)$] if
    \[
      \norm{f_n - f}_{\L^1(\mu)} =
      \int \abs{f_n - f} \dmu \taking{n \to \infty} 0.
    \]
  \item[(3) $f_n \taking{\mu} f$ in measure] if
    \[
      \forall \epsilon > 0 \quad
      \mu\del{\set{x \in X \colon \abs{f_n(x) - f(x)} \geq \epsilon}}
      \taking{n \to \infty} 0.
    \]
\end{description}
\begin{remark}
  We have that $(2) \Rightarrow (3)$ by Markov's inequality.
\end{remark}

\begin{remark}
  We have that $(3) \Rightarrow (1)$ for some subsequence.
  There exists a subsequence $\set{n_j}$ such that
  $f_{n_j} \taking{n \to \infty} f$ $\mu$-almost everywhere.
\end{remark}
\begin{remark}
  We have that $(1) \nRightarrow (2)$ and $(2) \nRightarrow (1)$.
\end{remark}

\begin{example}
  There are examples here.
\end{example}

%% Some examples talking

\subsection{Egoroff's theorem}
\begin{remark}
  The following theorem tells us when $(1) \Rightarrow (0)$.
\end{remark}
\begin{theorem}[Egoroff's theorem]
  Let $\mu(X) < \infty$ and
  $f_n \taking{n \to \infty} f$ $\mu$-almost everywhere.
  Then for all $\epsilon > 0$ there exists $E \subseteq X$ such
  that $\mu(E) < \epsilon$ and
  $f_n \taking{n \to \infty} f$ uniformly on $E^c$.
\end{theorem}
\begin{proof}
  To be added
\end{proof}

\newpage

\section{Approximations of functions in \texorpdfstring{$\L^1(\mu)$}{L}}

\begin{definition}[Support]
  Let $f \colon X \to \R$ be a function.
  The set-theoretic support of $f$ is
  \[
    \supp(f) = \set{x \in X \mid f(x) \neq 0}.
  \]
  The (topological) support of $f$ is
  \[
    \supp(f) = \overline{f^{-1}\del{\set{0}^c}}.
  \]
\end{definition}

\begin{proposition}
  The set of simple functions is dense in $\L^1(\mu)$.
  That is,
  we can choose $\phi = \sum_{j=1}^{k} c_j \mathbbm{1}_{E_j}$ such
  that $\int |f - \phi| \dmu < \epsilon$.
  We have $\phi \in \L^1(\mu)$ because
  $\norm{\phi} \le \underbrace{\norm{f}}_{\text{finite}} +
  \underbrace{\norm{\phi - f}}_{\epsilon} < \infty$.
\end{proposition}

\begin{proposition}
  Let $\mu$ be the Lebesgue--Stieltjes measure on $\R$.
  Then
  \begin{enumerate}
    \item[(1)] We can choose $E_j$ to be a finite disjoint union of open
      intervals for $\phi$.
    \item[(2)] The set of continuous functions $f \colon \R \to \R$
      with compact support (also denoted $C_c(\R)$) is dense in $\L^1(\mu)$.
  \end{enumerate}
\end{proposition}

\begin{proposition}
  Let $\mu = m$ in $\R^n$ (or any other Borel regular measure).
  Then
  \begin{enumerate}
    \item[(1)] We can choose $E_j$ to be a disjoint union of open boxes
      for $\phi$.
    \item[(2)] The set $C_c(\R^n)$
      (continuous functions $f \colon \R^n \to \R$ with compact support)
      is dense in $\L^1(m)$.
  \end{enumerate}
\end{proposition}
\begin{proof}
  To be added
\end{proof}

\begin{remark}
  The following theorem shows us that a measurable function
  is ``almost continuous''.
\end{remark}
\begin{theorem}[Lusin's theorem]
  Let $f \colon [a,b] \to \R$ be Lebesgue measurable.
  Then for all $\epsilon > 0$ there exists a compact $E \subseteq [a,b]$
  such that $m\del{[a,b] \setminus E} \le \epsilon$,
  and $f\vert_E$ is continuous.
\end{theorem}
\begin{remark}
  Being continuous on a subset of the domain $E \subset [a,b]$
  means that for any $\set{x_n}_{n=1}^{\infty} \subset E$ such that
  $x_n \to x'$ then $\lim_{n\to\infty} f(x_n) = f(x')$.
\end{remark}
\begin{proof}
  In the problems section.
\end{proof}

% Riemann integral reminder

\begin{theorem}
  Let $f \colon [a,b] \to \R$ be bounded.
  \begin{enumerate}
    \item[(1)] If $f$ is Riemann integrable, then $f$ is Lebesgue measurable
      (and thus also Lebesgue integrable),
      and $\int_{a}^{b} f(x) \dx = \int_{[a,b]} f \dm$.
    \item[(2)] $f$ is Riemann integrable if and only if
      $m\del{\set{x \in [a,b] \mid \text{$f$ has a discontinuity at $x$}}}
      = 0$.
  \end{enumerate}
\end{theorem}
\begin{proof}
  To be added.
\end{proof}
\begin{remark}
  This theorem is not true on a noncompact interval.
\end{remark}
\begin{example}
  Define the function
  \[
    f(x) :=
    \sum_{n=1}^{\floor{x}} \frac{(-1)^n}{n} \cdot \mathbbm{1}_{[n,n+1)}.
  \]
  Since $f$ is Lebesgue integrable if and only if $\abs{f}$ is Lebesgue
  integrable and the harmonic series diverges,
  we get that $f \notin \L^1(m)$.
  However, we have that
  \begin{align*}
    \int_{1}^{\infty} f(x) \dx &=
    \lim_{R \to \infty} \int_{1}^{\infty} f(x) \dx \\ &=
    \lim_{R \to R} \sum_{n=1}^{\floor{R/2}}
    \del{\frac{1}{2n} - \frac{1}{2n + 1}} \sim
    \sum_{n=1}^{\infty} \frac{1}{n^2} < \infty.
  \end{align*}
\end{example}
\begin{remark}
  This theorem is not true on an unbounded function.
\end{remark}
\begin{example}
  Define the function
  \[
    f(x) := \sum_{n=1}^{\floor{x}} (-1)^n \cdot n \cdot
    \mathbbm{1}_{\left[\frac{1}{n+1},\frac{1}{n}\right]}.
  \]
  We have that
  \[
    \int_{0}^{1} f(x) \dx =
    \lim_{\epsilon \to 0} \int_{\epsilon}^{1} f(x) \dx < \infty
  \]
  where the last inequality follows from conditional convergence.
  However, $\abs{f}$ is not Lebesgue integrable since the integral
  diverges for it.
\end{example}

% Push forward ----- meaning?

\newpage

\section{Product measures}
\begin{definition}[Product $\sigma$-algebra]
  Let $(X,\A_X,\mu_X)$ and $(Y,\A_Y,\mu_Y)$ be measure spaces.
  We define their product measure $\A_X \otimes \A_Y$ to be
  \[
    \A_X \otimes \A_Y :=
    \sigma\del{A \times B \mid A \in \A_X \stand B \in \A_Y}.
  \]
\end{definition}
We would like to construct a product measure $\mu_X \otimes \mu_Y$
on the product measurable space $(X \times Y, \A_X \otimes \A_Y)$
such that
\[
  \mu_X \otimes \mu_Y(A \times B) =
  \mu_X(A) \cdot \mu_Y(B).
\]
For example, if we denote Lebesgure measure on $\R^n$ as $m^n$,
we would like to have $m^1 \otimes m^1 = m^2$ and
$\bigotimes_{i=1}^{n} m^1 = m^n$ etc.
\begin{exercise}
  Show that $m^2(A \times B) = m^1(A) \cdot m^1(B)$.
  % by cases
\end{exercise}

Before we do that, we need a couple more definitions.

\begin{definition}[Semialgebra]
  A set $S \subseteq P(X)$ is a semialgebra over $X$ if
  \begin{enumerate}
    \item[(1)] $\emptyset \in \mathcal S$;
    \item[(2)] $E,F \in \mathcal E \implies E \cap F \in \mathcal E$;
    \item[(3)] $E \in \mathcal E \implies E^c$ is a finite disjoint
      union of elements in $\mathcal E$.
  \end{enumerate}
\end{definition}

\begin{definition}[Measurable box]
  A measurable box is a set of the form $A \times B \subseteq X \times Y$
  such that $A \in \A_X$ and $B \in \A_Y$.
  We denote this family with $\mathcal E$.
\end{definition}
\begin{remark}
  The set of measurable boxes $\mathcal E \subseteq P(X \times Y)$ is
  a semialgebra over $X \times Y$.
\end{remark}

We saw in the recitations that a fintite disjoint union of elements
of a semialgebra forms an algebra, thus a finite disjoint union
of measurable boxes forms an algebra. 
For this section we now set $\A = \sigma(\mathcal E)$.
To be added clarification

\begin{definition}[The $\pi$ function]
  We define the function $\pi \colon \mathcal E \to [0,\infty]$ with
  \[
    \pi\del{A \times B} := \mu_X(A) \times \mu_Y(B).
  \]
\end{definition}
\begin{proposition}
  Let $A \times B = \bigsqcup_{i=1}^{\infty} A_i \times B_i$.
  Then $\pi(A \times B) = \sum_{i=1}^{\infty} \pi(A_i \times B_i)$.
\end{proposition}
\begin{proof}
  First we notice that
  \[
    \ind_A(x) \cdot \ind_B(y) =
    \ind_{A \times B}(x,y) =
    \sum_{i=1}^{\infty} \ind_{A_i \times B_i}(x,y) =
    \sum_{i=1}^{\infty} \ind_{A_i}(x) \cdot \ind_{B_i}(y).
  \]
  If we consider $y$ as a parameter and integrate by $\dmu_X$ we get
  \begin{align*}
    \mu_X(A) \cdot \underbrace{\ind_B(y)}_{\text{constant}} &=
    \int \ind_A(x) \cdot \ind_B(y) \dmu_X \underset{\text{MCT}}{=}
    \sum_{i=1}^{\infty} \int \ind_{A_i}(x) \cdot \ind_{B_i}(y) \dmu_X \\ &=
    \sum_{i=1}^{\infty} \mu_X(A_i) \cdot \ind_{B_i}(y).
  \end{align*}
  We can now integrate by $\dmu_Y$ and get
  \[
    \mu_X(A) \cdot \mu_Y(B) =
    \sum_{i=1}^{\infty} \mu_X(A_i) \cdot \mu_Y(B_i)
  \]
  which completes the proof.
\end{proof}
\begin{remark}
  This implies that $\pi$ is $\sigma$-additive and in particular additive
  on $\mathcal E$.
\end{remark}

\begin{definition}[The extended $\pi$ function]
  We define the function $\pi \colon \A \to [0,\infty]$ with
  \[
    \pi\del{\bigsqcup_{i=1}^{N} A_i \times B_i} :=
    \sum_{i=1}^{N} \pi\del{A_i \times B_i}.
  \]
\end{definition}
\begin{proposition}
  This is a good definition.
\end{proposition}
\begin{proof}
  To be added.
\end{proof}

\begin{corollary}
  The function $\pi \colon A \to [0,\infty]$ is a pre-measure.
\end{corollary}
\begin{proof}
  To be added.
\end{proof}

According to Carath\'eodory's extension theorem
(\Cref{thm:car-extension-theorem}),
$\pi$ induces an outer measure $(\mu_X \otimes \mu_Y)^*$ on $X \times Y$,
so that its restriction to $\sigma(\A) = \A_X \otimes \A_Y$,
which we will denote with $\mu_X \otimes \mu_Y$ is a measure that
extends $\pi$.

More preciesly we have that for all $E \subseteq X \times Y$ that
\[
  (\mu_X \otimes \mu_Y)^*(E) :=
  \inf\set{\sum_{i=1}^{\infty} \mu_X(A_i) \cdot \mu_Y(B_i) \middle\vert
  E \subseteq \bigcup_{i=1}^{\infty} A_i \times B_i \st
  A_i \in \A_X \stand B_i \in \A_Y}.
\]




\begin{definition}[Bounded variation]
\end{definition}









\end{document}
